% Vietnamese translation for OI Public Library
% Source-PDF: 模板 TEMPLATES/计算几何模板-HIT.pdf
\documentclass[11pt]{article}
\usepackage{oipl-vn}
\setmonofont{WenQuanYi Zen Hei Mono}[Scale=MatchLowercase]

\title{Mẫu hình học tính toán}
\author{ACM@HIT\\jerrybond}
\date{}

\begin{document}
\maketitle
\tableofcontents

\section{Công thức hình học}

\subsection{Tam giác}
\begin{enumerate}
\item Nửa chu vi: $P=(a+b+c)/2$.
\item Diện tích: $S=aH_a/2=ab\sin C/2=\sqrt{P(P-a)(P-b)(P-c)}$.
\item Đường trung tuyến: $M_a=\sqrt{2(b^2+c^2)-a^2}/2=\sqrt{b^2+c^2+2bc\cos A}/2$.
\item Đường phân giác: $T_a=\sqrt{bc((b+c)^2-a^2)}/(b+c)=2bc\cos(A/2)/(b+c)$.
\item Đường cao: $H_a=b\sin C=c\sin B=\sqrt{b^2-\left((a^2+b^2-c^2)/(2a)\right)^2}$.
\item Bán kính đường tròn nội tiếp:
\[
\begin{aligned}
 r&=\frac{S}{P}
 =\frac{a\sin(B/2)\sin(C/2)}{\sin((B+C)/2)}\\
 &=4R\sin(A/2)\sin(B/2)\sin(C/2)
 =\sqrt{\frac{(P-a)(P-b)(P-c)}{P}}\\
 &=P\tan(A/2)\tan(B/2)\tan(C/2).
\end{aligned}
\]
\item Bán kính đường tròn ngoại tiếp: $R=abc/(4S)=a/(2\sin A)=b/(2\sin B)=c/(2\sin C)$.
\end{enumerate}

\subsection{Tứ giác}
Gọi $D_1,D_2$ là hai đường chéo, $M$ là đoạn nối trung điểm hai đường chéo, và $A$ là góc giữa hai đường chéo.
\begin{enumerate}
\item $a^2+b^2+c^2+d^2=D_1^2+D_2^2+4M^2$.
\item $S=D_1D_2\sin A/2$.
\item Với tứ giác nội tiếp đường tròn: $ac+bd=D_1D_2$.
\item Với tứ giác nội tiếp đường tròn: $S=\sqrt{(P-a)(P-b)(P-c)(P-d)}$, trong đó $P$ là nửa chu vi.
\end{enumerate}

\subsection{Đa giác đều n cạnh}
Gọi $R$ là bán kính đường tròn ngoại tiếp và $r$ là bán kính đường tròn nội tiếp.
\begin{enumerate}
\item Góc ở tâm: $A=2\pi/n$.
\item Góc trong: $C=(n-2)\pi/n$.
\item Độ dài cạnh: $a=2\sqrt{R^2-r^2}=2R\sin(A/2)=2r\tan(A/2)$.
\item Diện tích: $S=nar/2=nr^2\tan(A/2)=nR^2\sin A/2=na^2/(4\tan(A/2))$.
\end{enumerate}

\subsection{Hình tròn}
\begin{enumerate}
\item Độ dài cung: $l=rA$.
\item Độ dài dây cung: $a=2\sqrt{2hr-h^2}=2r\sin(A/2)$.
\item Chiều cao cung tròn: $h=r-\sqrt{r^2-a^2/4}=r(1-\cos(A/2))=a\tan(A/4)/2$.
\item Diện tích hình quạt: $S_1=rl/2=r^2A/2$.
\item Diện tích hình viên phân: $S_2=(rl-a(r-h))/2=r^2(A-\sin A)/2$.
\end{enumerate}

\subsection{Lăng trụ}
\begin{enumerate}
\item Thể tích: $V=Ah$, trong đó $A$ là diện tích đáy và $h$ là chiều cao.
\item Diện tích xung quanh: $S=lp$, trong đó $l$ là độ dài cạnh bên và $p$ là chu vi thiết diện vuông góc.
\item Diện tích toàn phần: $T=S+2A$.
\end{enumerate}

\subsection{Hình chóp}
\begin{enumerate}
\item Thể tích: $V=Ah/3$, trong đó $A$ là diện tích đáy và $h$ là chiều cao.
\item Với hình chóp đều, diện tích xung quanh: $S=lp/2$, trong đó $l$ là đường cao mặt bên và $p$ là chu vi đáy.
\item Với hình chóp đều, diện tích toàn phần: $T=S+A$.
\end{enumerate}

\subsection{Hình chóp cụt}
\begin{enumerate}
\item Thể tích: $V=(A_1+A_2+\sqrt{A_1A_2})h/3$, trong đó $A_1,A_2$ là diện tích hai đáy và $h$ là chiều cao.
\item Với chóp cụt đều, diện tích xung quanh: $S=(p_1+p_2)l/2$, trong đó $p_1,p_2$ là chu vi hai đáy và $l$ là đường cao mặt bên.
\item Diện tích toàn phần: $T=S+A_1+A_2$.
\end{enumerate}

\subsection{Hình trụ}
\begin{enumerate}
\item Diện tích xung quanh: $S=2\pi rh$.
\item Diện tích toàn phần: $T=2\pi r(h+r)$.
\item Thể tích: $V=\pi r^2h$.
\end{enumerate}

\subsection{Hình nón}
\begin{enumerate}
\item Đường sinh: $l=\sqrt{h^2+r^2}$.
\item Diện tích xung quanh: $S=\pi rl$.
\item Diện tích toàn phần: $T=\pi r(l+r)$.
\item Thể tích: $V=\pi r^2h/3$.
\end{enumerate}

\subsection{Hình nón cụt}
\begin{enumerate}
\item Đường sinh: $l=\sqrt{h^2+(r_1-r_2)^2}$.
\item Diện tích xung quanh: $S=\pi(r_1+r_2)l$.
\item Diện tích toàn phần: $T=\pi r_1(l+r_1)+\pi r_2(l+r_2)$.
\item Thể tích: $V=\pi(r_1^2+r_2^2+r_1r_2)h/3$.
\end{enumerate}

\subsection{Hình cầu}
\begin{enumerate}
\item Diện tích toàn phần: $T=4\pi r^2$.
\item Thể tích: $V=4\pi r^3/3$.
\end{enumerate}

\subsection{Đới cầu}
\begin{enumerate}
\item Diện tích xung quanh: $S=2\pi rh$.
\item Diện tích toàn phần: $T=\pi(2rh+r_1^2+r_2^2)$.
\item Thể tích: $V=\pi h(3(r_1^2+r_2^2)+h^2)/6$.
\end{enumerate}

\subsection{Quạt cầu}
\begin{enumerate}
\item Diện tích toàn phần: $T=\pi r(2h+r_0)$, trong đó $h$ là chiều cao chỏm cầu và $r_0$ là bán kính đáy chỏm cầu.
\item Thể tích: $V=2\pi r^2h/3$.
\end{enumerate}
\section{Đường thẳng và đoạn thẳng}
\subsection{Hàm chuẩn bị}
\begingroup
\scriptsize
\begin{verbatim}
//结构定义与宏定义
#include<stdio.h>
#include<string.h>
#include<stdlib.h>
#include <math.h>
#define eps 1e-8
#define zero(x) (((x)>0?(x):-(x))<eps)
struct point
{
   double x,y;
};
struct line
{
   point a,b;
};

//计算 cross product (P1-P0)x(P2-P0)
double xmult(point p1,point p2,point p0)
{
  return (p1.x-p0.x)*(p2.y-p0.y)-(p2.x-p0.x)*(p1.y-p0.y);
}
double xmult(double x1,double y1,double x2,double y2,double x0,double y0)
{
  return (x1-x0)*(y2-y0)-(x2-x0)*(y1-y0);
}

//计算 dot product (P1-P0).(P2-P0)
double dmult(point p1,point p2,point p0)
{
  return (p1.x-p0.x)*(p2.x-p0.x)+(p1.y-p0.y)*(p2.y-p0.y);
}
double dmult(double x1,double y1,double x2,double y2,double x0,double y0)
{
  return (x1-x0)*(x2-x0)+(y1-y0)*(y2-y0);




}

//两点距离
double distance(point p1,point p2)
{
  return sqrt((p1.x-p2.x)*(p1.x-p2.x)+(p1.y-p2.y)*(p1.y-p2.y));
}
double distance(double x1,double y1,double x2,double y2)
{
  return sqrt((x1-x2)*(x1-x2)+(y1-y2)*(y1-y2));
}


\end{verbatim}
\endgroup
\subsection{Kiểm tra ba điểm có thẳng hàng hay không}
\begingroup
\scriptsize
\begin{verbatim}
int dots_inline(point p1,point p2,point p3)
{
   return zero(xmult(p1,p2,p3));
}




\end{verbatim}
\endgroup
\subsection{Kiểm tra điểm có nằm trên đoạn thẳng hay không}
\begingroup
\scriptsize
\begin{verbatim}
//判点是否在线段上,包括端点（下面为两种接口模式）
int dot_online_in(point p,line l)
{
   return zero(xmult(p,l.a,l.b))&&(l.a.x-p.x)*(l.b.x-p.x)<eps&&(l.a.y-p.y)*(l.b.y-p.y)<eps;
}
int dot_online_in(point p,point l1,point l2)
{
   return zero(xmult(p,l1,l2))&&(l1.x-p.x)*(l2.x-p.x)<eps&&(l1.y-p.y)*(l2.y-p.y)<eps;
}
//判点是否在线段上,不包括端点
int dot_online_ex(point p,line l)
{
   return dot_online_in(p,l)&&(!zero(p.x-l.a.x)||!zero(p.y-l.a.y))
        &&(!zero(p.x-l.b.x)||!zero(p.y-l.b.y));
}




\end{verbatim}
\endgroup
\subsection{Kiểm tra hai điểm có nằm cùng phía của đoạn thẳng hay không}
\begingroup
\scriptsize
\begin{verbatim}
//判两点在线段同侧,点在线段上返回 0
int same_side(point p1,point p2,line l)
{
   return xmult(l.a,p1,l.b)*xmult(l.a,p2,l.b)>eps;
}
int same_side(point p1,point p2,point l1,point l2)
{
   return xmult(l1,p1,l2)*xmult(l1,p2,l2)>eps;
}


\end{verbatim}
\endgroup
\subsection{Kiểm tra hai điểm có nằm khác phía của đoạn thẳng hay không}
\begingroup
\scriptsize
\begin{verbatim}
//判两点在线段异侧,点在线段上返回 0
int opposite_side(point p1,point p2,line l)
{
   return xmult(l.a,p1,l.b)*xmult(l.a,p2,l.b)<-eps;
}
int opposite_side(point p1,point p2,point l1,point l2)
{
   return xmult(l1,p1,l2)*xmult(l1,p2,l2)<-eps;
}


\end{verbatim}
\endgroup
\subsection{Tìm điểm đối xứng qua đường thẳng}
\begingroup
\scriptsize
\begin{verbatim}
// 点关于直线的对称点 // by lyt
// 缺点：用了斜率
// 也可以利用"点到直线上的最近点"来做，避免使用斜率。
point symmetric_point(point p1, point l1, point l2)
{
   point ret;
   if (l1.x > l2.x - eps && l1.x < l2.x + eps)
   {
      ret.x = (2 * l1.x - p1.x);
      ret.y = p1.y;
   }



    else
    {
       double k = (l1.y - l2.y ) / (l1.x - l2.x);
       ret.x = (2*k*k*l1.x + 2*k*p1.y - 2*k*l1.y - k*k*p1.x + p1.x) / (1 + k*k);
       ret.y = p1.y - (ret.x - p1.x ) / k;
    }
    return ret;
}




\end{verbatim}
\endgroup
\subsection{Kiểm tra hai đoạn thẳng có giao nhau hay không}
\subsubsection{Phiên bản thường dùng}
\begingroup
\scriptsize
\begin{verbatim}
//定义点
struct Point
{
   double x;
   double y;
};
typedef struct Point point;

//叉积
double multi(point p0, point p1, point p2)
{
  return ( p1.x - p0.x )*( p2.y - p0.y )-( p2.x - p0.x )*( p1.y - p0.y );
}


//相交返回 true,否则为 false, 接口为两线段的端点
bool isIntersected(point s1,point e1, point s2,point e2)
{
  return (max(s1.x,e1.x) >= min(s2.x,e2.x)) &&
        (max(s2.x,e2.x) >= min(s1.x,e1.x)) &&
        (max(s1.y,e1.y) >= min(s2.y,e2.y)) &&
        (max(s2.y,e2.y) >= min(s1.y,e1.y)) &&
        (multi(s1,s2,e1)*multi(s1,e1,e2)>0) &&
        (multi(s2,s1,e2)*multi(s2,e2,e1)>0);
}




\end{verbatim}
\endgroup
\subsubsection{Phiên bản ít dùng}
\begingroup
\scriptsize
\begin{verbatim}
//判两线段相交,包括端点和部分重合
int intersect_in(line u,line v)
{
   if (!dots_inline(u.a,u.b,v.a)||!dots_inline(u.a,u.b,v.b))
      return !same_side(u.a,u.b,v)&&!same_side(v.a,v.b,u);
   return dot_online_in(u.a,v)||dot_online_in(u.b,v)||dot_online_in(v.a,u)||dot_online_in(v.b,u);
}
int intersect_in(point u1,point u2,point v1,point v2)
{
   if (!dots_inline(u1,u2,v1)||!dots_inline(u1,u2,v2))
      return !same_side(u1,u2,v1,v2)&&!same_side(v1,v2,u1,u2);
               return dot_online_in(u1,v1,v2)||dot_online_in(u2,v1,v2)||dot_online_in(v1,u1,u2)||
dot_online_in(v2,u1,u2);
}

//判两线段相交,不包括端点和部分重合
int intersect_ex(line u,line v)
{
   return opposite_side(u.a,u.b,v)&&opposite_side(v.a,v.b,u);
}
int intersect_ex(point u1,point u2,point v1,point v2)
{
   return opposite_side(u1,u2,v1,v2)&&opposite_side(v1,v2,u1,u2);
}


\end{verbatim}
\endgroup
\subsection{Tìm giao điểm của hai đường thẳng}
\begingroup
\scriptsize
\begin{verbatim}
//计算两直线交点,注意事先判断直线是否平行!
//线段交点请另外判线段相交(同时还是要判断是否平行!)
point intersection(point u1,point u2,point v1,point v2)
{
  point ret=u1;
  double t=((u1.x-v1.x)*(v1.y-v2.y)-(u1.y-v1.y)*(v1.x-v2.x))
         /((u1.x-u2.x)*(v1.y-v2.y)-(u1.y-u2.y)*(v1.x-v2.x));
  ret.x+=(u2.x-u1.x)*t;
  ret.y+=(u2.y-u1.y)*t;
  return ret;
}



\end{verbatim}
\endgroup
\subsection{Điểm gần nhất từ một điểm đến đường thẳng}
\begingroup
\scriptsize
\begin{verbatim}
point ptoline(point p,point l1,point l2)
{
  point t=p;
  t.x+=l1.y-l2.y,t.y+=l2.x-l1.x;
  return intersection(p,t,l1,l2);
}


\end{verbatim}
\endgroup
\subsection{Điểm gần nhất từ một điểm đến đoạn thẳng}
\begingroup
\scriptsize
\begin{verbatim}
point ptoseg(point p,point l1,point l2)
{
  point t=p;
  t.x+=l1.y-l2.y,t.y+=l2.x-l1.x;
  if (xmult(l1,t,p)*xmult(l2,t,p)>eps)
     return distance(p,l1)<distance(p,l2)?l1:l2;
  return intersection(p,t,l1,l2);
}



\end{verbatim}
\endgroup
\section{Đa giác}
\subsection{Hàm số thực chuẩn bị}
\begingroup
\scriptsize
\begin{verbatim}
#include <stdlib.h>
#include<stdio.h>
#include<string.h>
#include <math.h>
#define MAXN 1000

//offset 为多变形坐标的最大绝对值
#define offset 10000
#define eps 1e-8

//浮点数判 0



#define zero(x) (((x)>0?(x):-(x))<eps)

//浮点数判断符
#define _sign(x) ((x)>eps?1:((x)<-eps?2:0))

//定义点
struct point
{
   double x,y;
}pt[MAXN ];

//定义线段
struct line
{
   point a,b;
};

//叉积
double xmult(point p1,point p2,point p0)
{
  return (p1.x-p0.x)*(p2.y-p0.y)-(p2.x-p0.x)*(p1.y-p0.y);
}




\end{verbatim}
\endgroup
\subsection{Kiểm tra đa giác lồi}
\begingroup
\scriptsize
\begin{verbatim}
//判定凸多边形,顶点按顺时针或逆时针给出,允许相邻边共线,是凸多边形返回 1，否则返回
0
int is_convex(int n,point* p)
{
   int i,s[3]={1,1,1};
   for (i=0;i<n&&s[1]|s[2];i++)
      s[_sign(xmult(p[(i+1)%n],p[(i+2)%n],p[i]))]=0;
   return s[1]|s[2];
}

//判凸行，顶点按顺时针或逆时针给出,不允许相邻边共线,是凸多边形返回 1，否则返回 0
int is_convex_v2(int n,point* p)
{
   int i,s[3]={1,1,1};
   for (i=0;i<n&&s[0]&&s[1]|s[2];i++)
      s[_sign(xmult(p[(i+1)%n],p[(i+2)%n],p[i]))]=0;



    return s[0]&&s[1]|s[2];
}




\end{verbatim}
\endgroup
\subsection{Kiểm tra điểm có nằm trong đa giác hay không}
\begingroup
\scriptsize
\begin{verbatim}
//判点在凸多边形内或多边形边上时返回 1，严格在凸多边形外返回 0
int inside_convex(point q,int n,point* p)
{
   int i,s[3]={1,1,1};
   for (i=0;i<n&&s[1]|s[2];i++)
      s[_sign(xmult(p[(i+1)%n],q,p[i]))]=0;
   return s[1]|s[2];
}

//判点严格在凸多边形内返回 1,在边上或者严格在外返回 0
int inside_convex_v2(point q,int n,point* p)
{
   int i,s[3]={1,1,1};
   for (i=0;i<n&&s[0]&&s[1]|s[2];i++)
      s[_sign(xmult(p[(i+1)%n],q,p[i]))]=0;
   return s[0]&&s[1]|s[2];
}

//判点在任意多边形内,顶点按顺时针或逆时针给出
//on_edge 表示点在多边形边上时的返回值, offset 为多边形坐标上限,严格在内返回 1，严格
在外返回 0
int inside_polygon(point q,int n,point* p,int on_edge=2)
{
   point q2;
   int i=0,count;
   while (i<n)
      for (count=i=0,q2.x=rand()+offset,q2.y=rand()+offset;i<n;i++)
      {
         if (zero(xmult(q,p[i],p[(i+1)%n]))&&(p[i].x-q.x)*(p[(i+1)%n].x-q.x)<eps
            &&(p[i].y-q.y)*(p[(i+1)%n].y-q.y)<eps)
            return on_edge;

        else if (zero(xmult(q,q2,p[i])))
           break;

        else if (xmult(q,p[i],q2)*xmult(q,p[(i+1)%n],q2)<-eps&&



           xmult(p[i],q,p[(i+1)%n])*xmult(p[i],q2,p[(i+1)%n])<-eps)
           count++;
       }
    return count&1;
}




\end{verbatim}
\endgroup
\subsection{Kiểm tra đoạn thẳng có nằm trong đa giác bất kỳ hay không}
\begingroup
\scriptsize
\begin{verbatim}
//预备函数
inline int opposite_side(point p1,point p2,point l1,point l2)
{
   return xmult(l1,p1,l2)*xmult(l1,p2,l2)<-eps;
}
inline int dot_online_in(point p,point l1,point l2)
{
   return zero(xmult(p,l1,l2))&&(l1.x-p.x)*(l2.x-p.x)<eps&&(l1.y-p.y)*(l2.y-p.y)<eps;
}

//判线段在任意多边形内,顶点按顺时针或逆时针给出,与边界相交返回 1
int inside_polygon(point l1,point l2,int n,point* p)
{
   point t[MAXN],tt;
   int i,j,k=0;
   if (!inside_polygon(l1,n,p)||!inside_polygon(l2,n,p))
      return 0;
   for (i=0;i<n;i++)
   {
      if (opposite_side(l1,l2,p[i],p[(i+1)%n])&&opposite_side(p[i],p[(i+1)%n],l1,l2))
          return 0;
      else if (dot_online_in(l1,p[i],p[(i+1)%n]))
          t[k++]=l1;
      else if (dot_online_in(l2,p[i],p[(i+1)%n]))
          t[k++]=l2;
      else if (dot_online_in(p[i],l1,l2))
          t[k++]=p[i];
   }
   for (i=0;i<k;i++)
      for (j=i+1;j<k;j++)
      {
          tt.x=(t[i].x+t[j].x)/2;
          tt.y=(t[i].y+t[j].y)/2;



         if (!inside_polygon(tt,n,p))
            return 0;
       }
    return 1;
}



\end{verbatim}
\endgroup
\section{Tam giác}
\subsection{Hàm chuẩn bị}
\begingroup
\scriptsize
\begin{verbatim}
#include <math.h>
#include <string.h>
#include <stdlib.h>
#include<stdio.h>
//定义点
struct point
{
   double x,y;
};
typedef struct point point;

//定义直线
struct line
{
   point a,b;
};
typedef struct line line;
//两点距离
double distance(point p1,point p2)
{
     return sqrt((p1.x-p2.x)*(p1.x-p2.x)+(p1.y-p2.y)*(p1.y-p2.y));
}
//两直线求交点
point intersection(line u,line v)
{
     point ret=u.a;
     double t=((u.a.x-v.a.x)*(v.a.y-v.b.y)-(u.a.y-v.a.y)*(v.a.x-v.b.x))
                /((u.a.x-u.b.x)*(v.a.y-v.b.y)-(u.a.y-u.b.y)*(v.a.x-v.b.x));




     ret.x+=(u.b.x-u.a.x)*t;
     ret.y+=(u.b.y-u.a.y)*t;
     return ret;
}




\end{verbatim}
\endgroup
\subsection{Tìm tâm ngoại tiếp tam giác}
\begingroup
\scriptsize
\begin{verbatim}
point circumcenter(point a,point b,point c)
{
     line u,v;
     u.a.x=(a.x+b.x)/2;
     u.a.y=(a.y+b.y)/2;
     u.b.x=u.a.x-a.y+b.y;
     u.b.y=u.a.y+a.x-b.x;
     v.a.x=(a.x+c.x)/2;
     v.a.y=(a.y+c.y)/2;
     v.b.x=v.a.x-a.y+c.y;
     v.b.y=v.a.y+a.x-c.x;
     return intersection(u,v);
}




\end{verbatim}
\endgroup
\subsection{Tìm tâm nội tiếp tam giác}
\begingroup
\scriptsize
\begin{verbatim}
point incenter(point a,point b,point c)
{
     line u,v;
     double m,n;
     u.a=a;
     m=atan2(b.y-a.y,b.x-a.x);
     n=atan2(c.y-a.y,c.x-a.x);
     u.b.x=u.a.x+cos((m+n)/2);
     u.b.y=u.a.y+sin((m+n)/2);
     v.a=b;
     m=atan2(a.y-b.y,a.x-b.x);
     n=atan2(c.y-b.y,c.x-b.x);
     v.b.x=v.a.x+cos((m+n)/2);
     v.b.y=v.a.y+sin((m+n)/2);




     return intersection(u,v);
}


\end{verbatim}
\endgroup
\subsection{Tìm trực tâm tam giác}
\begingroup
\scriptsize
\begin{verbatim}
point perpencenter(point a,point b,point c)
{
     line u,v;
     u.a=c;
     u.b.x=u.a.x-a.y+b.y;
     u.b.y=u.a.y+a.x-b.x;
     v.a=b;
     v.b.x=v.a.x-a.y+c.y;
     v.b.y=v.a.y+a.x-c.x;
     return intersection(u,v);
}



\end{verbatim}
\endgroup
\section{Đường tròn}
\subsection{Hàm chuẩn bị}
\begingroup
\scriptsize
\begin{verbatim}
#include <math.h>
#include <stdlib.h>
#include <stdio.h>
#include <string.h>
#define eps 1e-8
struct point
{
   double x,y;
};
typedef struct point point;
double xmult(point p1,point p2,point p0)
{
   return (p1.x-p0.x)*(p2.y-p0.y)-(p2.x-p0.x)*(p1.y-p0.y);
}
double distance(point p1,point p2)
{



  return sqrt((p1.x-p2.x)*(p1.x-p2.x)+(p1.y-p2.y)*(p1.y-p2.y));
}
//点到直线的距离
double disptoline(point p,point l1,point l2)
{
  return fabs(xmult(p,l1,l2))/distance(l1,l2);
}
//求两直线交点
point intersection(point u1,point u2,point v1,point v2)
{
  point ret=u1;
  double t=((u1.x-v1.x)*(v1.y-v2.y)-(u1.y-v1.y)*(v1.x-v2.x))
         /((u1.x-u2.x)*(v1.y-v2.y)-(u1.y-u2.y)*(v1.x-v2.x));
  ret.x+=(u2.x-u1.x)*t;
  ret.y+=(u2.y-u1.y)*t;
  return ret;
}


\end{verbatim}
\endgroup
\subsection{Kiểm tra đường thẳng có cắt đường tròn hay không}
\begingroup
\scriptsize
\begin{verbatim}
//判直线和圆相交,包括相切
int intersect_line_circle(point c,double r,point l1,point l2)
{
   return disptoline(c,l1,l2)<r+eps;
}




\end{verbatim}
\endgroup
\subsection{Kiểm tra đoạn thẳng có cắt đường tròn hay không}
\begingroup
\scriptsize
\begin{verbatim}
int intersect_seg_circle(point c,double r, point l1,point l2)
{
   double t1=distance(c,l1)-r,t2=distance(c,l2)-r;
   point t=c;
   if (t1<eps||t2<eps)
      return t1>-eps||t2>-eps;
   t.x+=l1.y-l2.y;
   t.y+=l2.x-l1.x;
   return xmult(l1,c,t)*xmult(l2,c,t)<eps&&disptoline(c,l1,l2)-r<eps;
}





\end{verbatim}
\endgroup
\subsection{Kiểm tra hai đường tròn có giao nhau hay không}
\begingroup
\scriptsize
\begin{verbatim}
int intersect_circle_circle(point c1,double r1,point c2,double r2)
{
   return distance(c1,c2)<r1+r2+eps&&distance(c1,c2)>fabs(r1-r2)-eps;
}


\end{verbatim}
\endgroup
\subsection{Tính điểm trên đường tròn gần điểm p nhất}
\begingroup
\scriptsize
\begin{verbatim}
//当 p 为圆心时，返回圆心本身
point dot_to_circle(point c,double r,point p)
{
  point u,v;
  if (distance(p,c)<eps)
     return p;
  u.x=c.x+r*fabs(c.x-p.x)/distance(c,p);
  u.y=c.y+r*fabs(c.y-p.y)/distance(c,p)*((c.x-p.x)*(c.y-p.y)<0?-1:1);
  v.x=c.x-r*fabs(c.x-p.x)/distance(c,p);
  v.y=c.y-r*fabs(c.y-p.y)/distance(c,p)*((c.x-p.x)*(c.y-p.y)<0?-1:1);
  return distance(u,p)<distance(v,p)?u:v;
}


\end{verbatim}
\endgroup
\subsection{Tính giao điểm của đường thẳng và đường tròn}
\begingroup
\scriptsize
\begin{verbatim}
//计算直线与圆的交点,保证直线与圆有交点
//计算线段与圆的交点可用这个函数后判点是否在线段上
void intersection_line_circle(point c,double r,point l1,point l2,point& p1,point& p2)
{
  point p=c;
  double t;
  p.x+=l1.y-l2.y;
  p.y+=l2.x-l1.x;
  p=intersection(p,c,l1,l2);
  t=sqrt(r*r-distance(p,c)*distance(p,c))/distance(l1,l2);
  p1.x=p.x+(l2.x-l1.x)*t;
  p1.y=p.y+(l2.y-l1.y)*t;
  p2.x=p.x-(l2.x-l1.x)*t;
  p2.y=p.y-(l2.y-l1.y)*t;
}



\end{verbatim}
\endgroup
\subsection{Tính giao điểm của hai đường tròn}
\begingroup
\scriptsize
\begin{verbatim}
//计算圆与圆的交点,保证圆与圆有交点,圆心不重合
void intersection_circle_circle(point c1,double r1,point c2,double r2,point& p1,point& p2)
{
  point u,v;
  double t;
  t=(1+(r1*r1-r2*r2)/distance(c1,c2)/distance(c1,c2))/2;
  u.x=c1.x+(c2.x-c1.x)*t;
  u.y=c1.y+(c2.y-c1.y)*t;
  v.x=u.x+c1.y-c2.y;
  v.y=u.y-c1.x+c2.x;
  intersection_line_circle(c1,r1,u,v,p1,p2);
}



\end{verbatim}
\endgroup
\section{Mặt cầu}
\subsection{Tính góc ở tâm từ kinh độ và vĩ độ trên Trái Đất}
\begingroup
\scriptsize
\begin{verbatim}
#include <math.h>
const double pi=acos(-1);

//计算圆心角 lat 表示纬度,-90<=w<=90,lng 表示经度
//返回两点所在大圆劣弧对应圆心角,0<=angle<=pi
double angle(double lng1,double lat1,double lng2,double lat2)
{
  double dlng=fabs(lng1-lng2)*pi/180;
  while (dlng>=pi+pi)
     dlng-=pi+pi;
  if (dlng>pi)
     dlng=pi+pi-dlng;
  lat1*=pi/180,lat2*=pi/180;
  return acos(cos(lat1)*cos(lat2)*cos(dlng)+sin(lat1)*sin(lat2));
}





\end{verbatim}
\endgroup
\subsection{Tính khoảng cách đường thẳng giữa hai điểm trên Trái Đất từ kinh độ, vĩ độ}
\begingroup
\scriptsize
\begin{verbatim}
//计算距离,r 为球半径
double line_dist(double r,double lng1,double lat1,double lng2,double lat2)
{
  double dlng=fabs(lng1-lng2)*pi/180;
  while (dlng>=pi+pi)
     dlng-=pi+pi;
  if (dlng>pi)
     dlng=pi+pi-dlng;
  lat1*=pi/180,lat2*=pi/180;
  return r*sqrt(2-2*(cos(lat1)*cos(lat2)*cos(dlng)+sin(lat1)*sin(lat2)));
}


\end{verbatim}
\endgroup
\subsection{Tính khoảng cách mặt cầu giữa hai điểm trên Trái Đất từ kinh độ, vĩ độ}
\begingroup
\scriptsize
\begin{verbatim}
//计算球面距离,r 为球半径
inline double sphere_dist(double r,double lng1,double lat1,double lng2,double lat2)
{
   return r*angle(lng1,lat1,lng2,lat2);
}



\end{verbatim}
\endgroup
\section{Một số mẫu hình học ba chiều}
\subsection{Hàm chuẩn bị}
\begingroup
\scriptsize
\begin{verbatim}
//三维几何函数库
#include <math.h>
#define eps 1e-8
#define zero(x) (((x)>0?(x):-(x))<eps)
struct point3{double x,y,z;};
struct line3{point3 a,b;};
struct plane3{point3 a,b,c;};

//计算 cross product U x V



point3 xmult(point3 u,point3 v){
     point3 ret;
     ret.x=u.y*v.z-v.y*u.z;
     ret.y=u.z*v.x-u.x*v.z;
     ret.z=u.x*v.y-u.y*v.x;
     return ret;
}

//计算 dot product U . V
double dmult(point3 u,point3 v){
    return u.x*v.x+u.y*v.y+u.z*v.z;
}

//矢量差 U - V
point3 subt(point3 u,point3 v){
     point3 ret;
     ret.x=u.x-v.x;
     ret.y=u.y-v.y;
     ret.z=u.z-v.z;
     return ret;
}

//取平面法向量
point3 pvec(plane3 s){
     return xmult(subt(s.a,s.b),subt(s.b,s.c));
}
point3 pvec(point3 s1,point3 s2,point3 s3){
     return xmult(subt(s1,s2),subt(s2,s3));
}

//两点距离,单参数取向量大小
double distance(point3 p1,point3 p2){
    return sqrt((p1.x-p2.x)*(p1.x-p2.x)+(p1.y-p2.y)*(p1.y-p2.y)+(p1.z-p2.z)*(p1.z-p2.z));
}

//向量大小
double vlen(point3 p){
    return sqrt(p.x*p.x+p.y*p.y+p.z*p.z);
}





\end{verbatim}
\endgroup
\subsection{Kiểm tra ba điểm có thẳng hàng hay không}
\begingroup
\scriptsize
\begin{verbatim}
//判三点共线
int dots_inline(point3 p1,point3 p2,point3 p3){
     return vlen(xmult(subt(p1,p2),subt(p2,p3)))<eps;
}




\end{verbatim}
\endgroup
\subsection{Kiểm tra bốn điểm có đồng phẳng hay không}
\begingroup
\scriptsize
\begin{verbatim}
//判四点共面
int dots_onplane(point3 a,point3 b,point3 c,point3 d){
     return zero(dmult(pvec(a,b,c),subt(d,a)));
}


\end{verbatim}
\endgroup
\subsection{Kiểm tra điểm có nằm trên đoạn thẳng hay không}
\begingroup
\scriptsize
\begin{verbatim}
//判点是否在线段上,包括端点和共线
int dot_online_in(point3 p,line3 l){
     return zero(vlen(xmult(subt(p,l.a),subt(p,l.b))))&&(l.a.x-p.x)*(l.b.x-p.x)<eps&&
          (l.a.y-p.y)*(l.b.y-p.y)<eps&&(l.a.z-p.z)*(l.b.z-p.z)<eps;
}
int dot_online_in(point3 p,point3 l1,point3 l2){
     return zero(vlen(xmult(subt(p,l1),subt(p,l2))))&&(l1.x-p.x)*(l2.x-p.x)<eps&&
          (l1.y-p.y)*(l2.y-p.y)<eps&&(l1.z-p.z)*(l2.z-p.z)<eps;
}

//判点是否在线段上,不包括端点
int dot_online_ex(point3 p,line3 l){
     return dot_online_in(p,l)&&(!zero(p.x-l.a.x)||!zero(p.y-l.a.y)||!zero(p.z-l.a.z))&&
          (!zero(p.x-l.b.x)||!zero(p.y-l.b.y)||!zero(p.z-l.b.z));
}
int dot_online_ex(point3 p,point3 l1,point3 l2){
     return dot_online_in(p,l1,l2)&&(!zero(p.x-l1.x)||!zero(p.y-l1.y)||!zero(p.z-l1.z))&&
          (!zero(p.x-l2.x)||!zero(p.y-l2.y)||!zero(p.z-l2.z));
}




\end{verbatim}
\endgroup
\subsection{Kiểm tra điểm có nằm trên tam giác trong không gian hay không}
\begingroup
\scriptsize
\begin{verbatim}
//判点是否在空间三角形上,包括边界,三点共线无意义
int dot_inplane_in(point3 p,plane3 s){
     return zero(vlen(xmult(subt(s.a,s.b),subt(s.a,s.c)))-vlen(xmult(subt(p,s.a),subt(p,s.b)))-
          vlen(xmult(subt(p,s.b),subt(p,s.c)))-vlen(xmult(subt(p,s.c),subt(p,s.a))));
}
int dot_inplane_in(point3 p,point3 s1,point3 s2,point3 s3){
     return zero(vlen(xmult(subt(s1,s2),subt(s1,s3)))-vlen(xmult(subt(p,s1),subt(p,s2)))-
          vlen(xmult(subt(p,s2),subt(p,s3)))-vlen(xmult(subt(p,s3),subt(p,s1))));
}



//判点是否在空间三角形上,不包括边界,三点共线无意义
int dot_inplane_ex(point3 p,plane3 s){
     return dot_inplane_in(p,s)&&vlen(xmult(subt(p,s.a),subt(p,s.b)))>eps&&
          vlen(xmult(subt(p,s.b),subt(p,s.c)))>eps&&vlen(xmult(subt(p,s.c),subt(p,s.a)))>eps;
}
int dot_inplane_ex(point3 p,point3 s1,point3 s2,point3 s3){
     return dot_inplane_in(p,s1,s2,s3)&&vlen(xmult(subt(p,s1),subt(p,s2)))>eps&&
          vlen(xmult(subt(p,s2),subt(p,s3)))>eps&&vlen(xmult(subt(p,s3),subt(p,s1)))>eps;
}


\end{verbatim}
\endgroup
\subsection{Kiểm tra hai điểm có nằm cùng phía của đoạn thẳng hay không}
\begingroup
\scriptsize
\begin{verbatim}
//判两点在线段同侧,点在线段上返回 0,不共面无意义
int same_side(point3 p1,point3 p2,line3 l){
      return dmult(xmult(subt(l.a,l.b),subt(p1,l.b)),xmult(subt(l.a,l.b),subt(p2,l.b)))>eps;
}
int same_side(point3 p1,point3 p2,point3 l1,point3 l2){
      return dmult(xmult(subt(l1,l2),subt(p1,l2)),xmult(subt(l1,l2),subt(p2,l2)))>eps;
}


\end{verbatim}
\endgroup
\subsection{Kiểm tra hai điểm có nằm khác phía của đoạn thẳng hay không}
\begingroup
\scriptsize
\begin{verbatim}
//判两点在线段异侧,点在线段上返回 0,不共面无意义
int opposite_side(point3 p1,point3 p2,line3 l){
     return dmult(xmult(subt(l.a,l.b),subt(p1,l.b)),xmult(subt(l.a,l.b),subt(p2,l.b)))<-eps;
}



int opposite_side(point3 p1,point3 p2,point3 l1,point3 l2){
     return dmult(xmult(subt(l1,l2),subt(p1,l2)),xmult(subt(l1,l2),subt(p2,l2)))<-eps;
}




\end{verbatim}
\endgroup
\subsection{Kiểm tra hai điểm có nằm cùng phía của mặt phẳng hay không}
\begingroup
\scriptsize
\begin{verbatim}
//判两点在平面同侧,点在平面上返回 0
int same_side(point3 p1,point3 p2,plane3 s){
      return dmult(pvec(s),subt(p1,s.a))*dmult(pvec(s),subt(p2,s.a))>eps;
}
int same_side(point3 p1,point3 p2,point3 s1,point3 s2,point3 s3){
      return dmult(pvec(s1,s2,s3),subt(p1,s1))*dmult(pvec(s1,s2,s3),subt(p2,s1))>eps;
}




\end{verbatim}
\endgroup
\subsection{Kiểm tra hai điểm có nằm khác phía của mặt phẳng hay không}
\begingroup
\scriptsize
\begin{verbatim}
//判两点在平面异侧,点在平面上返回 0
int opposite_side(point3 p1,point3 p2,plane3 s){
     return dmult(pvec(s),subt(p1,s.a))*dmult(pvec(s),subt(p2,s.a))<-eps;
}
int opposite_side(point3 p1,point3 p2,point3 s1,point3 s2,point3 s3){
     return dmult(pvec(s1,s2,s3),subt(p1,s1))*dmult(pvec(s1,s2,s3),subt(p2,s1))<-eps;
}


\end{verbatim}
\endgroup
\subsection{Kiểm tra hai đường thẳng trong không gian có song song hay không}
\begingroup
\scriptsize
\begin{verbatim}
//判两直线平行
int parallel(line3 u,line3 v){
     return vlen(xmult(subt(u.a,u.b),subt(v.a,v.b)))<eps;
}
int parallel(point3 u1,point3 u2,point3 v1,point3 v2){
     return vlen(xmult(subt(u1,u2),subt(v1,v2)))<eps;
}





\end{verbatim}
\endgroup
\subsection{Kiểm tra hai mặt phẳng có song song hay không}
\begingroup
\scriptsize
\begin{verbatim}
//判两平面平行
int parallel(plane3 u,plane3 v){
     return vlen(xmult(pvec(u),pvec(v)))<eps;
}
int parallel(point3 u1,point3 u2,point3 u3,point3 v1,point3 v2,point3 v3){
     return vlen(xmult(pvec(u1,u2,u3),pvec(v1,v2,v3)))<eps;
}


\end{verbatim}
\endgroup
\subsection{Kiểm tra đường thẳng có song song với mặt phẳng hay không}
\begingroup
\scriptsize
\begin{verbatim}
//判直线与平面平行
int parallel(line3 l,plane3 s){
     return zero(dmult(subt(l.a,l.b),pvec(s)));
}
int parallel(point3 l1,point3 l2,point3 s1,point3 s2,point3 s3){
     return zero(dmult(subt(l1,l2),pvec(s1,s2,s3)));
}




\end{verbatim}
\endgroup
\subsection{Kiểm tra hai đường thẳng có vuông góc hay không}
\begingroup
\scriptsize
\begin{verbatim}
//判两直线垂直
int perpendicular(line3 u,line3 v){
     return zero(dmult(subt(u.a,u.b),subt(v.a,v.b)));
}
int perpendicular(point3 u1,point3 u2,point3 v1,point3 v2){
     return zero(dmult(subt(u1,u2),subt(v1,v2)));
}


\end{verbatim}
\endgroup
\subsection{Kiểm tra hai mặt phẳng có vuông góc hay không}
\begingroup
\scriptsize
\begin{verbatim}
//判两平面垂直
int perpendicular(plane3 u,plane3 v){



     return zero(dmult(pvec(u),pvec(v)));
}
int perpendicular(point3 u1,point3 u2,point3 u3,point3 v1,point3 v2,point3 v3){
     return zero(dmult(pvec(u1,u2,u3),pvec(v1,v2,v3)));
}


\end{verbatim}
\endgroup
\subsection{Kiểm tra hai đoạn thẳng trong không gian có giao nhau hay không}
\begingroup
\scriptsize
\begin{verbatim}
//判两线段相交,包括端点和部分重合
int intersect_in(line3 u,line3 v){
      if (!dots_onplane(u.a,u.b,v.a,v.b))
            return 0;
      if (!dots_inline(u.a,u.b,v.a)||!dots_inline(u.a,u.b,v.b))
            return !same_side(u.a,u.b,v)&&!same_side(v.a,v.b,u);
      return dot_online_in(u.a,v)||dot_online_in(u.b,v)||dot_online_in(v.a,u)||dot_online_in(v.b,u);
}
int intersect_in(point3 u1,point3 u2,point3 v1,point3 v2){
      if (!dots_onplane(u1,u2,v1,v2))
            return 0;
      if (!dots_inline(u1,u2,v1)||!dots_inline(u1,u2,v2))
            return !same_side(u1,u2,v1,v2)&&!same_side(v1,v2,u1,u2);
      return             dot_online_in(u1,v1,v2)||dot_online_in(u2,v1,v2)||dot_online_in(v1,u1,u2)||
dot_online_in(v2,u1,u2);
}

//判两线段相交,不包括端点和部分重合
int intersect_ex(line3 u,line3 v){
      return dots_onplane(u.a,u.b,v.a,v.b)&&opposite_side(u.a,u.b,v)&&opposite_side(v.a,v.b,u);
}
int intersect_ex(point3 u1,point3 u2,point3 v1,point3 v2){
      return
dots_onplane(u1,u2,v1,v2)&&opposite_side(u1,u2,v1,v2)&&opposite_side(v1,v2,u1,u2);
}


\end{verbatim}
\endgroup
\subsection{Kiểm tra đoạn thẳng có giao với tam giác trong không gian hay không}
\begingroup
\scriptsize
\begin{verbatim}
//判线段与空间三角形相交,包括交于边界和(部分)包含
int intersect_in(line3 l,plane3 s){
      return !same_side(l.a,l.b,s)&&!same_side(s.a,s.b,l.a,l.b,s.c)&&
           !same_side(s.b,s.c,l.a,l.b,s.a)&&!same_side(s.c,s.a,l.a,l.b,s.b);




}
int intersect_in(point3 l1,point3 l2,point3 s1,point3 s2,point3 s3){
      return !same_side(l1,l2,s1,s2,s3)&&!same_side(s1,s2,l1,l2,s3)&&
           !same_side(s2,s3,l1,l2,s1)&&!same_side(s3,s1,l1,l2,s2);
}

//判线段与空间三角形相交,不包括交于边界和(部分)包含
int intersect_ex(line3 l,plane3 s){
      return opposite_side(l.a,l.b,s)&&opposite_side(s.a,s.b,l.a,l.b,s.c)&&
           opposite_side(s.b,s.c,l.a,l.b,s.a)&&opposite_side(s.c,s.a,l.a,l.b,s.b);
}
int intersect_ex(point3 l1,point3 l2,point3 s1,point3 s2,point3 s3){
      return opposite_side(l1,l2,s1,s2,s3)&&opposite_side(s1,s2,l1,l2,s3)&&
           opposite_side(s2,s3,l1,l2,s1)&&opposite_side(s3,s1,l1,l2,s2);
}


\end{verbatim}
\endgroup
\subsection{Tính giao điểm của hai đường thẳng}
\begingroup
\scriptsize
\begin{verbatim}
//计算两直线交点,注意事先判断直线是否共面和平行!
//线段交点请另外判线段相交(同时还是要判断是否平行!)
point3 intersection(line3 u,line3 v){
     point3 ret=u.a;
     double t=((u.a.x-v.a.x)*(v.a.y-v.b.y)-(u.a.y-v.a.y)*(v.a.x-v.b.x))
                 /((u.a.x-u.b.x)*(v.a.y-v.b.y)-(u.a.y-u.b.y)*(v.a.x-v.b.x));
     ret.x+=(u.b.x-u.a.x)*t;
     ret.y+=(u.b.y-u.a.y)*t;
     ret.z+=(u.b.z-u.a.z)*t;
     return ret;
}
point3 intersection(point3 u1,point3 u2,point3 v1,point3 v2){
     point3 ret=u1;
     double t=((u1.x-v1.x)*(v1.y-v2.y)-(u1.y-v1.y)*(v1.x-v2.x))
                 /((u1.x-u2.x)*(v1.y-v2.y)-(u1.y-u2.y)*(v1.x-v2.x));
     ret.x+=(u2.x-u1.x)*t;
     ret.y+=(u2.y-u1.y)*t;
     ret.z+=(u2.z-u1.z)*t;
     return ret;
}





\end{verbatim}
\endgroup
\subsection{Tính giao điểm của đường thẳng và mặt phẳng}
\begingroup
\scriptsize
\begin{verbatim}
//计算直线与平面交点,注意事先判断是否平行,并保证三点不共线!
//线段和空间三角形交点请另外判断
point3 intersection(line3 l,plane3 s){
     point3 ret=pvec(s);
     double t=(ret.x*(s.a.x-l.a.x)+ret.y*(s.a.y-l.a.y)+ret.z*(s.a.z-l.a.z))/
          (ret.x*(l.b.x-l.a.x)+ret.y*(l.b.y-l.a.y)+ret.z*(l.b.z-l.a.z));
     ret.x=l.a.x+(l.b.x-l.a.x)*t;
     ret.y=l.a.y+(l.b.y-l.a.y)*t;
     ret.z=l.a.z+(l.b.z-l.a.z)*t;
     return ret;
}
point3 intersection(point3 l1,point3 l2,point3 s1,point3 s2,point3 s3){
     point3 ret=pvec(s1,s2,s3);
     double t=(ret.x*(s1.x-l1.x)+ret.y*(s1.y-l1.y)+ret.z*(s1.z-l1.z))/
          (ret.x*(l2.x-l1.x)+ret.y*(l2.y-l1.y)+ret.z*(l2.z-l1.z));
     ret.x=l1.x+(l2.x-l1.x)*t;
     ret.y=l1.y+(l2.y-l1.y)*t;
     ret.z=l1.z+(l2.z-l1.z)*t;
     return ret;
}


\end{verbatim}
\endgroup
\subsection{Tính giao tuyến của hai mặt phẳng}
\begingroup
\scriptsize
\begin{verbatim}
//计算两平面交线,注意事先判断是否平行,并保证三点不共线!
line3 intersection(plane3 u,plane3 v){
     line3 ret;
     ret.a=parallel(v.a,v.b,u.a,u.b,u.c)?
intersection(v.b,v.c,u.a,u.b,u.c):intersection(v.a,v.b,u.a,u.b,u.c);
     ret.b=parallel(v.c,v.a,u.a,u.b,u.c)?
intersection(v.b,v.c,u.a,u.b,u.c):intersection(v.c,v.a,u.a,u.b,u.c);
     return ret;
}
line3 intersection(point3 u1,point3 u2,point3 u3,point3 v1,point3 v2,point3 v3){
     line3 ret;
     ret.a=parallel(v1,v2,u1,u2,u3)?intersection(v2,v3,u1,u2,u3):intersection(v1,v2,u1,u2,u3);
     ret.b=parallel(v3,v1,u1,u2,u3)?intersection(v2,v3,u1,u2,u3):intersection(v3,v1,u1,u2,u3);
     return ret;
}



\end{verbatim}
\endgroup
\subsection{Khoảng cách từ điểm đến đường thẳng}
\begingroup
\scriptsize
\begin{verbatim}
//点到直线距离
double ptoline(point3 p,line3 l){
    return vlen(xmult(subt(p,l.a),subt(l.b,l.a)))/distance(l.a,l.b);
}
double ptoline(point3 p,point3 l1,point3 l2){
    return vlen(xmult(subt(p,l1),subt(l2,l1)))/distance(l1,l2);
}


\end{verbatim}
\endgroup
\subsection{Tính khoảng cách từ điểm đến mặt phẳng}
\begingroup
\scriptsize
\begin{verbatim}
//点到平面距离
double ptoplane(point3 p,plane3 s){
    return fabs(dmult(pvec(s),subt(p,s.a)))/vlen(pvec(s));
}
double ptoplane(point3 p,point3 s1,point3 s2,point3 s3){
    return fabs(dmult(pvec(s1,s2,s3),subt(p,s1)))/vlen(pvec(s1,s2,s3));
}


\end{verbatim}
\endgroup
\subsection{Tính khoảng cách giữa hai đường thẳng}
\begingroup
\scriptsize
\begin{verbatim}
//直线到直线距离
double linetoline(line3 u,line3 v){
    point3 n=xmult(subt(u.a,u.b),subt(v.a,v.b));
    return fabs(dmult(subt(u.a,v.a),n))/vlen(n);
}
double linetoline(point3 u1,point3 u2,point3 v1,point3 v2){
    point3 n=xmult(subt(u1,u2),subt(v1,v2));
    return fabs(dmult(subt(u1,v1),n))/vlen(n);
}




\end{verbatim}
\endgroup
\subsection{Giá trị cos của góc giữa hai đường thẳng trong không gian}
\begingroup
\scriptsize
\begin{verbatim}
//两直线夹角 cos 值
double angle_cos(line3 u,line3 v){



     return dmult(subt(u.a,u.b),subt(v.a,v.b))/vlen(subt(u.a,u.b))/vlen(subt(v.a,v.b));
}
double angle_cos(point3 u1,point3 u2,point3 v1,point3 v2){
    return dmult(subt(u1,u2),subt(v1,v2))/vlen(subt(u1,u2))/vlen(subt(v1,v2));
}


\end{verbatim}
\endgroup
\subsection{Giá trị cos của góc giữa hai mặt phẳng}
\begingroup
\scriptsize
\begin{verbatim}
//两平面夹角 cos 值
double angle_cos(plane3 u,plane3 v){
    return dmult(pvec(u),pvec(v))/vlen(pvec(u))/vlen(pvec(v));
}
double angle_cos(point3 u1,point3 u2,point3 u3,point3 v1,point3 v2,point3 v3){
    return dmult(pvec(u1,u2,u3),pvec(v1,v2,v3))/vlen(pvec(u1,u2,u3))/vlen(pvec(v1,v2,v3));
}


\end{verbatim}
\endgroup
\subsection{Giá trị sin của góc giữa đường thẳng và mặt phẳng}
\begingroup
\scriptsize
\begin{verbatim}
//直线平面夹角 sin 值
double angle_sin(line3 l,plane3 s){
    return dmult(subt(l.a,l.b),pvec(s))/vlen(subt(l.a,l.b))/vlen(pvec(s));
}
double angle_sin(point3 l1,point3 l2,point3 s1,point3 s2,point3 s3){
    return dmult(subt(l1,l2),pvec(s1,s2,s3))/vlen(subt(l1,l2))/vlen(pvec(s1,s2,s3));
}



\end{verbatim}
\endgroup
\section{Khoảng cách Manhattan xa nhất}
\begingroup
\scriptsize
\begin{verbatim}
#include <stdio.h>
#define INF 9999999999999.0
struct Point
{
   double x[5];
}pt[100005];
double dis[32][100005], coe[5], minx[32], maxx[32];
//去掉绝对值后有 2^D 种可能
void GetD(int N, int D)



{
    int s, i, j, tot=(1<<D);
    for (s=0;s<tot;s++)
    {
       for (i=0;i<D;i++)
          if (s&(1<<i))
              coe[i]=-1.0;
          else coe[i]=1.0;
       for (i=0;i<N;i++)
       {
          dis[s][i]=0.0;
          for (j=0;j<D;j++)
              dis[s][i]=dis[s][i]+coe[j]*pt[i].x[j];
       }
    }
}
//取每种可能中的最大差距
void Solve(int N, int D)
{
   int s, i, tot=(1<<D);
   double tmp, ans;
   for (s=0;s<tot;s++)
   {
      minx[s]=INF;
      maxx[s]=-INF;
      for (i=0; i<N; i++)
      {
         if (minx[s]>dis[s][i]) minx[s]=dis[s][i];
         if (maxx[s]<dis[s][i]) maxx[s]=dis[s][i];
      }
   }
   ans=0.0;
   for (s=0; s<tot; s++)
   {
      tmp=maxx[s]-minx[s];
      if (tmp>ans) ans=tmp;
   }
   printf("%.2lf\n", ans);
}
int main (void)
{
   int n, i;
   while (scanf("%d",&n)==1)
   {


      for (i=0;i<n;i++)
       scanf("%lf%lf%lf%lf%lf",&pt[i].x[0],&pt[i].x[1],&pt[i].x[2],&pt[i].x[3],&pt[i].x[4]);
      GetD(n, 5);
      Solve(n, 5);
    }
    return 0;
}



\end{verbatim}
\endgroup
\section{Cặp điểm gần nhất}
\begingroup
\scriptsize
\begin{verbatim}
#include <stdio.h>
#include <math.h>
#include <stdlib.h>
#define Max(x,y) (x)>(y)?(x):(y)
struct Q
{
   double x, y;
}q[100001], sl[10], sr[10];

int cntl, cntr, lm, rm;
double ans;

int cmp(const void*p1, const void*p2)
{
   struct Q*a1=(struct Q*)p1;
   struct Q*a2=(struct Q*)p2;
   if (a1->x<a2->x)return -1;
   else if (a1->x==a2->x)return 0;
   else return 1;
}

double CalDis(double x1, double y1, double x2, double y2)
{
  return sqrt((x1-x2)*(x1-x2)+(y1-y2)*(y1-y2));
}

void MinDis(int l, int r)
{
  if (l==r) return;
  double dis;



    if (l+1==r)
    {
       dis=CalDis(q[l].x,q[l].y,q[r].x,q[r].y);
       if (ans>dis) ans=dis;
       return;
    }
    int mid=(l+r)>>1, i, j;
    MinDis(l,mid);
    MinDis(mid+1,r);

    lm=mid+1-5;
    if (lm<l) lm=l;
    rm=mid+5;
    if (rm>r) rm=r;

    cntl=cntr=0;
    for (i=mid;i>=lm;i--)
    {
      if (q[mid+1].x-q[i].x>=ans)break;
      sl[++cntl]=q[i];
    }
    for (i=mid+1;i<=rm;i++)
    {
      if (q[i].x-q[mid].x>=ans)break;
      sr[++cntr]=q[i];
    }

    for (i=1;i<=cntl;i++)
      for (j=1;j<=cntr;j++)
      {
          dis=CalDis(sl[i].x,sl[i].y,sr[j].x,sr[j].y);
          if (dis<ans) ans=dis;
      }
}

int main (void)
{
   int n, i;
   while (scanf("%d",&n)==1&&n)
   {
      for (i=1;i<=n;i++)
         scanf("%lf %lf", &q[i].x,&q[i].y);
      qsort(q+1,n,sizeof(struct Q),cmp);
      ans=CalDis(q[1].x,q[1].y,q[2].x,q[2].y);


      MinDis(1,n);
      printf("%.2lf\n",ans/2.0);
    }
    return 0;
}



\end{verbatim}
\endgroup
\section{Cặp điểm gần nhất}
\begingroup
\scriptsize
\begin{verbatim}
#include <stdio.h>
#include <math.h>
#include <stdlib.h>
#define Max(x,y) (x)>(y)?(x):(y)
struct Q
{
   double x, y;
}q[100001], sl[10], sr[10];

int cntl, cntr, lm, rm;
double ans;

int cmp(const void*p1, const void*p2)
{
   struct Q*a1=(struct Q*)p1;
   struct Q*a2=(struct Q*)p2;
   if (a1->x<a2->x)return -1;
   else if (a1->x==a2->x)return 0;
   else return 1;
}

double CalDis(double x1, double y1, double x2, double y2)
{
  return sqrt((x1-x2)*(x1-x2)+(y1-y2)*(y1-y2));
}

void MinDis(int l, int r)
{
  if (l==r) return;
  double dis;
  if (l+1==r)
  {



      dis=CalDis(q[l].x,q[l].y,q[r].x,q[r].y);
      if (ans>dis) ans=dis;
      return;
    }
    int mid=(l+r)>>1, i, j;
    MinDis(l,mid);
    MinDis(mid+1,r);

    lm=mid+1-5;
    if (lm<l) lm=l;
    rm=mid+5;
    if (rm>r) rm=r;

    cntl=cntr=0;
    for (i=mid;i>=lm;i--)
    {
      if (q[mid+1].x-q[i].x>=ans)break;
      sl[++cntl]=q[i];
    }
    for (i=mid+1;i<=rm;i++)
    {
      if (q[i].x-q[mid].x>=ans)break;
      sr[++cntr]=q[i];
    }

    for (i=1;i<=cntl;i++)
      for (j=1;j<=cntr;j++)
      {
          dis=CalDis(sl[i].x,sl[i].y,sr[j].x,sr[j].y);
          if (dis<ans) ans=dis;
      }
}

int main (void)
{
   int n, i;
   while (scanf("%d",&n)==1&&n)
   {
      for (i=1;i<=n;i++)
         scanf("%lf %lf", &q[i].x,&q[i].y);
      qsort(q+1,n,sizeof(struct Q),cmp);
      ans=CalDis(q[1].x,q[1].y,q[2].x,q[2].y);
      MinDis(1,n);
      printf("%.2lf\n",ans/2.0);


    }
    return 0;
}



\end{verbatim}
\endgroup
\section{Đường tròn bao nhỏ nhất}
\begingroup
\scriptsize
\begin{verbatim}
#include<stdio.h>
#include<string.h>
#include<math.h>
struct Point
{
   double x;
   double y;
}pt[1005];
struct Traingle
{
   struct Point p[3];
};
struct Circle
{
   struct Point center;
   double r;
}ans;
//计算两点距离
double Dis(struct Point p, struct Point q)
{
   double dx=p.x-q.x;
   double dy=p.y-q.y;
   return sqrt(dx*dx+dy*dy);
}
//计算三角形面积
double Area(struct Traingle ct)
{
     return fabs((ct.p[1].x-ct.p[0].x)*(ct.p[2].y-ct.p[0].y)-(ct.p[2].x-ct.p[0].x)*(ct.p[1].y-ct.p[0].y))/
2.0;
}
//求三角形的外接圆，返回圆心和半径(存在结构体"圆"中)
struct Circle CircumCircle(struct Traingle t)
{
   struct Circle tmp;



  double a, b, c, c1, c2;
  double xA, yA, xB, yB, xC, yC;
  a = Dis(t.p[0], t.p[1]);
  b = Dis(t.p[1], t.p[2]);
  c = Dis(t.p[2], t.p[0]);
  //根据 S = a * b * c / R / 4;求半径 R
  tmp.r = (a*b*c)/(Area(t)*4.0);
  xA = t.p[0].x;
  yA = t.p[0].y;
  xB = t.p[1].x;
  yB = t.p[1].y;
  xC = t.p[2].x;
  yC = t.p[2].y;
  c1 = (xA*xA+yA*yA - xB*xB-yB*yB) / 2;
  c2 = (xA*xA+yA*yA - xC*xC-yC*yC) / 2;
  tmp.center.x = (c1*(yA - yC)-c2*(yA - yB)) / ((xA - xB)*(yA - yC)-(xA - xC)*(yA - yB));
  tmp.center.y = (c1*(xA - xC)-c2*(xA - xB)) / ((yA - yB)*(xA - xC)-(yA - yC)*(xA - xB));
  return tmp;
}
//确定最小包围圆
struct Circle MinCircle(int num, struct Traingle ct)
{
   struct Circle ret;
   if (num==0) ret.r = 0.0;
   else if (num==1)
   {
      ret.center = ct.p[0];
      ret.r = 0.0;
   }
   else if (num==2)
   {
      ret.center.x = (ct.p[0].x+ct.p[1].x)/2.0;
      ret.center.y = (ct.p[0].y+ct.p[1].y)/2.0;
      ret.r = Dis(ct.p[0], ct.p[1])/2.0;
   }
   else if(num==3) ret = CircumCircle(ct);
   return ret;
}
//递归实现增量算法
void Dfs(int x, int num, struct Traingle ct)
{
   int i, j;
   struct Point tmp;
   ans = MinCircle(num, ct);


  if (num==3) return;
  for (i=1; i<=x; i++)
     if (Dis(pt[i], ans.center)>ans.r)
     {
        ct.p[num]=pt[i];
        Dfs(i-1, num+1, ct);
        tmp=pt[i];
        for (j=i;j>=2;j--)
           pt[j]=pt[j-1];
        pt[1]=tmp;
     }
}
void Solve(int n)
{
   struct Traingle ct;
   Dfs(n, 0, ct);
}
int main (void)
{
   int n, i;
   while (scanf("%d", &n)!=EOF && n)
   {
      for (i=1;i<=n;i++)
         scanf("%lf %lf", &pt[i].x, &pt[i].y);
      Solve(n);
      printf("%.2lf %.2lf %.2lf\n", ans.center.x, ans.center.y, ans.r);
   }
   return 0;
}




\end{verbatim}
\endgroup
\section{Tìm giao điểm của hai đường tròn}
\begingroup
\scriptsize
\begin{verbatim}
#include<stdio.h>
#include<string.h>
#include<math.h>
#include<stdlib.h>
const double eps = 1e-8;
const double PI = acos(-1.0);





struct Point
{
   double x;
   double y;
};
typedef struct Point point;

struct Line
{
   double s, t;
};
typedef struct Line Line;

struct Circle
{
   Point center;
   double r;
   Line line[505];
   int cnt;
   bool covered;

}circle[105];

double distance(point p1, point p2)
{
  double dx = p1.x-p2.x;
  double dy = p1.y-p2.y;
  return sqrt(dx*dx + dy*dy);
}

point intersection(point u1,point u2, point v1,point v2)
{
  point ret = u1;
  double t=((u1.x-v1.x)*(v1.y-v2.y)-(u1.y-v1.y)*(v1.x-v2.x)) /
         ((u1.x-u2.x)*(v1.y-v2.y)-(u1.y-u2.y)*(v1.x-v2.x));
  ret.x += (u2.x-u1.x)*t;
  ret.y += (u2.y-u1.y)*t;
  return ret;
}

void intersection_line_circle(point c,double r,point l1,point l2,point& p1,point& p2)
{
  point p=c;
  double t;


    p.x+=l1.y-l2.y;
    p.y+=l2.x-l1.x;
    p=intersection(p,c,l1,l2);
    t=sqrt(r*r-distance(p,c)*distance(p,c))/distance(l1,l2);
    p1.x=p.x+(l2.x-l1.x)*t;
    p1.y=p.y+(l2.y-l1.y)*t;
    p2.x=p.x-(l2.x-l1.x)*t;
    p2.y=p.y-(l2.y-l1.y)*t;
}

//计算圆与圆的交点,保证圆与圆有交点,圆心不重合
void intersection_circle_circle(point c1,double r1,point c2,double r2,point& p1,point& p2)
{
  point u,v;
  double t;
  t=(1+(r1*r1-r2*r2)/distance(c1,c2)/distance(c1,c2))/2;
  u.x=c1.x+(c2.x-c1.x)*t;
  u.y=c1.y+(c2.y-c1.y)*t;
  v.x=u.x+c1.y-c2.y;
  v.y=u.y-c1.x+c2.x;
  intersection_line_circle(c1,r1,u,v,p1,p2);
}



\end{verbatim}
\endgroup
\section{Tìm tâm đường tròn ngoại tiếp tam giác}
\begingroup
\scriptsize
\begin{verbatim}
struct Point
{
   double x;
   double y;
}pt[1005];
struct Traingle
{
   struct Point p[3];
};
struct Circle
{
   struct Point center;
   double r;
}ans;
//计算两点距离



double Dis(struct Point p, struct Point q)
{
   double dx=p.x-q.x;
   double dy=p.y-q.y;
   return sqrt(dx*dx+dy*dy);
}
//计算三角形面积
double Area(struct Traingle ct)
{
     return fabs((ct.p[1].x-ct.p[0].x)*(ct.p[2].y-ct.p[0].y)-(ct.p[2].x-ct.p[0].x)*(ct.p[1].y-ct.p[0].y))/
2.0;
}
//求三角形的外接圆，返回圆心和半径(存在结构体"圆"中)
struct Circle CircumCircle(struct Traingle t)
{
   struct Circle tmp;
   double a, b, c, c1, c2;
   double xA, yA, xB, yB, xC, yC;
   a = Dis(t.p[0], t.p[1]);
   b = Dis(t.p[1], t.p[2]);
   c = Dis(t.p[2], t.p[0]);
   //根据 S = a * b * c / R / 4;求半径 R
   tmp.r = (a*b*c)/(Area(t)*4.0);
   xA = t.p[0].x;
   yA = t.p[0].y;
   xB = t.p[1].x;
   yB = t.p[1].y;
   xC = t.p[2].x;
   yC = t.p[2].y;
   c1 = (xA*xA+yA*yA - xB*xB-yB*yB) / 2;
   c2 = (xA*xA+yA*yA - xC*xC-yC*yC) / 2;
   tmp.center.x = (c1*(yA - yC)-c2*(yA - yB)) / ((xA - xB)*(yA - yC)-(xA - xC)*(yA - yB));
   tmp.center.y = (c1*(xA - xC)-c2*(xA - xB)) / ((yA - yB)*(xA - xC)-(yA - yC)*(xA - xB));
   return tmp;
}





\end{verbatim}
\endgroup
\section{Tìm bao lồi}
\begingroup
\scriptsize
\begin{verbatim}
#include <stdio.h>
#include <string.h>
#include <stdlib.h>
#include <math.h>
#define INF 999999999.9
#define PI acos(-1.0)
struct Point
{
   double x, y, dis;
}pt[1005], stack[1005], p0;
int top, tot;
//计算几何距离
double Dis(double x1, double y1, double x2, double y2)
{
   return sqrt((x1-x2)*(x1-x2)+(y1-y2)*(y1-y2));
}
//极角比较， 返回-1: p0p1 在 p0p2 的右侧，返回 0:p0,p1,p2 共线
int Cmp_PolarAngel(struct Point p1, struct Point p2, struct Point pb)
{
   double delta=(p1.x-pb.x)*(p2.y-pb.y)-(p2.x-pb.x)*(p1.y-pb.y);
   if (delta<0.0) return 1;
   else if (delta==0.0) return 0;
   else return -1;
}
// 判断向量 p2p3 是否对 p1p2 构成左旋
bool Is_LeftTurn(struct Point p3, struct Point p2, struct Point p1)
{
   int type=Cmp_PolarAngel(p3, p1, p2);
   if (type<0) return true;
   return false;
}
//先按极角排，再按距离由小到大排
int Cmp(const void*p1, const void*p2)
{
   struct Point*a1=(struct Point*)p1;
   struct Point*a2=(struct Point*)p2;
   int type=Cmp_PolarAngel(*a1, *a2, p0);
   if (type<0) return -1;



  else if (type==0)
  {
     if (a1->dis<a2->dis) return -1;
     else if (a1->dis==a2->dis) return 0;
     else return 1;
  }
  else return 1;
}
//求凸包
void Solve(int n)
{
  int i, k;
  p0.x=p0.y=INF;
  for (i=0;i<n;i++)
  {
     scanf("%lf %lf",&pt[i].x, &pt[i].y);
     if (pt[i].y < p0.y)
     {
         p0.y=pt[i].y;
         p0.x=pt[i].x;
         k=i;
     }
     else if (pt[i].y==p0.y)
     {
         if (pt[i].x<p0.x)
         {
            p0.x=pt[i].x;
            k=i;
         }
     }
  }
  pt[k]=pt[0];
  pt[0]=p0;
  for (i=1;i<n;i++)
     pt[i].dis=Dis(pt[i].x,pt[i].y, p0.x,p0.y);
  qsort(pt+1, n-1, sizeof(struct Point), Cmp);
  //去掉极角相同的点
  tot=1;
  for (i=2;i<n;i++)
     if (Cmp_PolarAngel(pt[i], pt[i-1], p0))
         pt[tot++]=pt[i-1];
  pt[tot++]=pt[n-1];
  //求凸包
  top=1;


  stack[0]=pt[0];
  stack[1]=pt[1];
  for (i=2;i<tot;i++)
  {
     while (top>=1 && Is_LeftTurn(pt[i], stack[top], stack[top-1])==false)
        top--;
     stack[++top]=pt[i];
  }
}
int main (void)
{
   int n;
   while (scanf("%d",&n)==2)
   {
      Solve(n);
   }
   return 0;
}



\end{verbatim}
\endgroup
\section{Bao lồi và thước cặp xoay để tìm mọi cặp đối cực, cặp điểm xa nhất}
\begingroup
\scriptsize
\begin{verbatim}
#include <stdio.h>
#include <string.h>
#include <stdlib.h>
#include <math.h>
#define INF 999999999.9
#define PI acos(-1.0)
struct Point
{
   double x, y, dis;
}pt[6005], stack[6005], p0;
int top, tot;
//计算几何距离
double Dis(double x1, double y1, double x2, double y2)
{
   return sqrt((x1-x2)*(x1-x2)+(y1-y2)*(y1-y2));
}
//极角比较， 返回-1: p0p1 在 p0p2 的右侧，返回 0:p0,p1,p2 共线
int Cmp_PolarAngel(struct Point p1, struct Point p2, struct Point pb)
{



  double delta=(p1.x-pb.x)*(p2.y-pb.y)-(p2.x-pb.x)*(p1.y-pb.y);
  if (delta<0.0) return 1;
  else if (delta==0.0) return 0;
  else return -1;
}
// 判断向量 p2p3 是否对 p1p2 构成左旋
bool Is_LeftTurn(struct Point p3, struct Point p2, struct Point p1)
{
   int type=Cmp_PolarAngel(p3, p1, p2);
   if (type<0) return true;
   return false;
}
//先按极角排，再按距离由小到大排
int Cmp(const void*p1, const void*p2)
{
   struct Point*a1=(struct Point*)p1;
   struct Point*a2=(struct Point*)p2;
   int type=Cmp_PolarAngel(*a1, *a2, p0);
   if (type<0) return -1;
   else if (type==0)
   {
      if (a1->dis<a2->dis) return -1;
      else if (a1->dis==a2->dis) return 0;
      else return 1;
   }
   else return 1;
}
//求凸包
void Hull(int n)
{
   int i, k;
   p0.x=p0.y=INF;
   for (i=0;i<n;i++)
   {
      scanf("%lf %lf",&pt[i].x, &pt[i].y);
      if (pt[i].y < p0.y)
      {
          p0.y=pt[i].y;
          p0.x=pt[i].x;
          k=i;
      }
      else if (pt[i].y==p0.y)
      {
          if (pt[i].x<p0.x)


        {
            p0.x=pt[i].x;
            k=i;
        }
    }
  }
  pt[k]=pt[0];
  pt[0]=p0;
  for (i=1;i<n;i++)
     pt[i].dis=Dis(pt[i].x,pt[i].y, p0.x,p0.y);
  qsort(pt+1, n-1, sizeof(struct Point), Cmp);
  //去掉极角相同的点
  tot=1;
  for (i=2;i<n;i++)
     if (Cmp_PolarAngel(pt[i], pt[i-1], p0))
        pt[tot++]=pt[i-1];
  pt[tot++]=pt[n-1];
  //求凸包
  top=1;
  stack[0]=pt[0];
  stack[1]=pt[1];
  for (i=2;i<tot;i++)
  {
     while (top>=1 && Is_LeftTurn(pt[i], stack[top], stack[top-1])==false)
        top--;
     stack[++top]=pt[i];
  }
}
//计算叉积
double CrossProduct(struct Point p1, struct Point p2, struct Point p3)
{
  return (p1.x-p3.x)*(p2.y-p3.y)-(p2.x-p3.x)*(p1.y-p3.y);
}
//卡壳旋转，求出凸多边形所有对踵点
void Rotate(struct Point*ch, int n)
{
  int i, p=1;
  double t1, t2, ans=0.0, dif;
  ch[n]=ch[0];
  for (i=0;i<n;i++)
  {
     //如果下一个点与当前边构成的三角形的面积更大，则说明此时不构成对踵点
     while (fabs(CrossProduct(ch[i],ch[i+1],ch[p+1])) > fabs(CrossProduct(ch[i],ch[i+1],ch[p])))
         p=(p+1)%n;


      dif=fabs(CrossProduct(ch[i],ch[i+1],ch[p+1])) - fabs(CrossProduct(ch[i],ch[i+1],ch[p]));
       //如果当前点和下一个点分别构成的三角形面积相等，则说明两条边即为平行线，对
角线两端都可能是对踵点
      if (dif==0.0)
      {
         t1=Dis(ch[p].x, ch[p].y, ch[i].x, ch[i].y);
         t2=Dis(ch[p+1].x, ch[p+1].y, ch[i+1].x, ch[i+1].y);
         if (t1>ans)ans=t1;
         if (t2>ans)ans=t2;
      }
      //说明 p，i 是对踵点
      else if (dif<0.0)
      {
         t1=Dis(ch[p].x, ch[p].y, ch[i].x, ch[i].y);
         if (t1>ans)ans=t1;
      }
   }
   printf("%.2lf\n",ans);
}
int main (void)
{
   int n;
   while (scanf("%d",&n)==1)
   {
      Hull(n);
      Rotate(stack, top+1);
   }
   return 0;
}




\end{verbatim}
\endgroup
\section{Bao lồi và thước cặp xoay để tìm tam giác có diện tích lớn nhất trên mặt phẳng}
\begingroup
\scriptsize
\begin{verbatim}
#include <stdio.h>
#include <string.h>
#include <stdlib.h>
#include <math.h>
#define INF 99999999999.9
#define PI acos(-1.0)
struct Point



{
   double x, y, dis;
}pt[50005], stack[50005], p0;
int top, tot;
double Dis(double x1, double y1, double x2, double y2)
{
   return sqrt((x1-x2)*(x1-x2)+(y1-y2)*(y1-y2));
}
int Cmp_PolarAngel(struct Point p1, struct Point p2, struct Point pb)
{
   double delta=(p1.x-pb.x)*(p2.y-pb.y)-(p2.x-pb.x)*(p1.y-pb.y);
   if (delta<0.0) return 1;
   else if (delta==0.0) return 0;
   else return -1;
}
bool Is_LeftTurn(struct Point p3, struct Point p2, struct Point p1)
{
   int type=Cmp_PolarAngel(p3, p1, p2);
   if (type<0) return true;
   return false;
}
int Cmp(const void*p1, const void*p2)
{
   struct Point*a1=(struct Point*)p1;
   struct Point*a2=(struct Point*)p2;
   int type=Cmp_PolarAngel(*a1, *a2, p0);
   if (type<0) return -1;
   else if (type==0)
   {
      if (a1->dis<a2->dis) return -1;
      else if (a1->dis==a2->dis) return 0;
      else return 1;
   }
   else return 1;
}
void Hull(int n)
{
   int i, k;
   p0.x=p0.y=INF;
   for (i=0;i<n;i++)
   {
      scanf("%lf %lf",&pt[i].x, &pt[i].y);
      if (pt[i].y < p0.y)
      {


       p0.y=pt[i].y;
       p0.x=pt[i].x;
       k=i;
    }
    else if (pt[i].y==p0.y)
    {
       if (pt[i].x<p0.x)
       {
          p0.x=pt[i].x;
          k=i;
       }
    }
  }
  pt[k]=pt[0];
  pt[0]=p0;
  for (i=1;i<n;i++)
     pt[i].dis=Dis(pt[i].x,pt[i].y, p0.x,p0.y);
  qsort(pt+1, n-1, sizeof(struct Point), Cmp);
  tot=1;
  for (i=2;i<n;i++)
     if (Cmp_PolarAngel(pt[i], pt[i-1], p0))
        pt[tot++]=pt[i-1];
  pt[tot++]=pt[n-1];
  top=1;
  stack[0]=pt[0];
  stack[1]=pt[1];
  for (i=2;i<tot;i++)
  {
     while (top>=1 && Is_LeftTurn(pt[i], stack[top], stack[top-1])==false)
        top--;
     stack[++top]=pt[i];
  }
}
double TArea(struct Point p1, struct Point p2, struct Point p3)
{
  return fabs((p1.x-p3.x)*(p2.y-p3.y)-(p2.x-p3.x)*(p1.y-p3.y));
}
void Rotate(struct Point*ch, int n)
{
  if (n<3)
  {
     printf("0.00\n");
     return;
  }


  int i, j, k;
  double ans=0.0, tmp;
  ch[n]=ch[0];
  for (i=0;i<n;i++)
  {
     j=(i+1)%n;
     k=(j+1)%n;
     while ((j!=k) && (k!=i))
     {
         while (TArea(ch[i],ch[j],ch[k+1])>TArea(ch[i],ch[j],ch[k]))
            k=(k+1)%n;
         tmp=TArea(ch[i],ch[j], ch[k]);
         if (tmp>ans) ans=tmp;
         j=(j+1)%n;
     }
  }
  printf("%.2lf\n",ans/2.0);
}
int main (void)
{
   int n;
   while (scanf("%d",&n)==1)
   {
      if (n==-1)break;
      Hull(n);
      Rotate(stack, top+1);
   }
   return 0;
}



\end{verbatim}
\endgroup
\section{Định lý Pick}
\begingroup
\scriptsize
\begin{verbatim}
// Pick 定理求整点多边形内部整点数目
// (1) 给定顶点座标均是整点（或正方形格点）的简单多边形，皮克定理说明了其面积 A 和
内部格点数目 i、边上格点数目 b 的关系：A = i + b/2 - 1；
// (2) 在两点（x1，y1），（ x2，y2）连线之间的整点个数（包含一个端点）为：gcd（|x1
－x2|，|y1－y2|）；
// (3) 求三角形面积用叉乘

#include<stdio.h>



#include<stdlib.h>
#include<math.h>
#include<string.h>

long long x[3], y[3], area, b;
long long My_Abs(long long t)
{
   if (t<0) return -t;
   return t;
}
long long Gcd(long long x, long long y)
{
   if (y==0) return x;
   long long mod=x%y;
   while (mod)
   {
      x=y;
      y=mod;
      mod=x%y;
   }
   return y;
}
int main (void)
{
   int i;
   while (1)
   {
      for (i = 0;i < 3;i ++)
          scanf("%lld %lld", &x[i], &y[i]);
      if(x[0]==0&&y[0]==0&&x[1]==0&&y[1]==0&&x[2]==0&&y[2]==0) break;
      area = (x[1]-x[0])*(y[2]-y[0])-(x[2]-x[0])*(y[1]-y[0]);
      area = My_Abs(area);
      b=0;
        b=Gcd(My_Abs(x[1]-x[0]), My_Abs(y[1]-y[0])) + Gcd(My_Abs(x[2]-x[0]), My_Abs(y[2]-
y[0])) + Gcd(My_Abs(x[1]-x[2]), My_Abs(y[1]-y[2]));
      printf("%lld\n", (area-b+2)/2);
   }
   return 0;
}





\end{verbatim}
\endgroup
\section{Tính diện tích và trọng tâm đa giác}
\begingroup
\scriptsize
\begin{verbatim}
#include <stdio.h>
#include <math.h>
int x[1000003], y[1000003];
double A, tx, ty, tmp;
int main (void)
{
   int cases, n, i;
   scanf ("%d", &cases);
   while (cases --)
   {
      scanf ("%d", &n);
      A = 0.0;
      x[0] = y[0] = 0;
      for (i = 1; i <= n; i ++)
      {
         scanf ("%d %d", &x[i], &y[i]);
         A += (x[i-1]*y[i] - x[i]*y[i-1]);
      }
      A += x[n]*y[1] - x[1]*y[n];
      A = A / 2.0;
      tx = ty = 0.0;
      for (i = 1; i < n; i ++)
      {
         tmp = x[i]*y[i+1] - x[i+1]*y[i];
         tx += (x[i]+x[i+1]) * tmp;
         ty += (y[i]+y[i+1]) * tmp;
      }
      tmp = x[n]*y[1] - x[1]*y[n];
      tx += (x[n]+x[1])*tmp;
      ty += (y[n]+y[1])*tmp;
      printf ("%.2lf %.2lf\n", tx/(6.0*A), ty/(6.0*A));
   }
   return 0;
}





\end{verbatim}
\endgroup
\section{Kiểm tra đa giác đơn có nhân hay không}
\begingroup
\scriptsize
\begin{verbatim}
#include <stdio.h>
#include <string.h>
const int INF = (1<<30);
struct Point
{
   int x, y;
}pt[150];
typedef struct Point Point;
bool turn_right[150];
int det(Point s1, Point t1, Point s2, Point t2)
{
   int d1x = t1.x-s1.x;
   int d1y = t1.y-s1.y;

  int d2x = t2.x-s2.x;
  int d2y = t2.y-s2.y;

  return d1x*d2y - d2x*d1y;
}
void Swap(int &a, int &b)
{
   if (a>b)
   {
      int t=a;
      a=b;
      b=t;
   }
}
int main (void)
{
   int n, i, cross, maxx, minx, maxy, miny, maxn, minn, countn=0;
   while (scanf("%d", &n)==1&&n)
   {
      maxx=maxy=-INF;
      minx=miny=INF;
      //点按顺时针给出
      for (i=1; i<=n; i++)
      {



     scanf("%d %d", &pt[i].x, &pt[i].y);
     if (maxx<pt[i].x) maxx=pt[i].x;
     if (maxy<pt[i].y) maxy=pt[i].y;
     if (minx>pt[i].x) minx=pt[i].x;
     if (miny>pt[i].y) miny=pt[i].y;
  }
  pt[n+1]=pt[1];
  pt[n+2]=pt[2];
  pt[n+3]=pt[3];
  pt[n+4]=pt[4];
  //求每条线段的转向
  for (i=1; i<=n+1; i ++)
  {
     cross = det(pt[i],pt[i+1], pt[i+1], pt[i+2]);
     if (cross<0)
        turn_right[i+1]=true;
     else turn_right[i+1]=false;
  }
  //两条边连续右转的为凸处，只有此时才可影响“核”肯恩存在的范围
  for (i=2; i<= n+1; i++)
     if (turn_right[i] && turn_right[i+1])
     {
        if (pt[i].x==pt[i+1].x)
        {
           minn=pt[i].y;
           maxn=pt[i+1].y;
           Swap(minn, maxn);
           if (minn>miny) miny=minn;
           if (maxn<maxy) maxy=maxn;
        }
        else
        {
           minn=pt[i].x;
           maxn=pt[i+1].x;
           Swap(minn, maxn);
           if (minn>minx) minx=minn;
           if (maxn<maxx) maxx=maxn;
        }
     }
  if (minx<=maxx && miny<=maxy)
     printf("Floor #%d\nSurveillance is possible.\n\n", ++countn);
  else printf("Floor #%d\nSurveillance is impossible.\n\n", ++countn);
}
return 0;


}




\end{verbatim}
\endgroup
\section{Mô phỏng luyện kim}
\begingroup
\scriptsize
\begin{verbatim}
#include <stdio.h>
#include <stdlib.h>
#include <math.h>
#define Lim 0.999999
#define EPS 1e-2
#define PI acos(-1.0)
double Temp, maxx, minx, maxy, miny, lx, ly, dif;
int nt, ns, nc;
struct Target
{
   double x, y;
}T[105];
struct Solution
{
   double x, y;
   double f;
}S[25], P, A;
double Dis(double x1, double y1, double x2, double y2)
{
   return sqrt((x1-x2)*(x1-x2)+(y1-y2)*(y1-y2));
}
void Seed(void)
{
   int i, j;
   for (i=0;i<ns;i++)
   {
      S[i].x=minx+((double)(rand()%1000+1)/1000.0)*lx;
      S[i].y=miny+((double)(rand()%1000+1)/1000.0)*ly;
      S[i].f=0.0;
      for (j=0;j<nt;j++)
          S[i].f=S[i].f+Dis(S[i].x,S[i].y, T[j].x, T[j].y);
   }
}
void Trans(void)
{



  int i, j, k;
  double theta;
  for (i=0;i<ns;i++)
  {
     P=S[i];
     for (j=0;j<nc;j++)
     {
         theta=(((double)(rand()%1000+1))/1000.0)*2.0*PI;
         A.x=P.x+Temp*cos(theta);
         A.y=P.y+Temp*sin(theta);
         if (A.x<minx||A.x>maxx||A.y<miny||A.y>maxy)
            continue;
         A.f=0.0;
         for (k=0;k<nt;k++)
            A.f=A.f+Dis(A.x,A.y,T[k].x,T[k].y);
         dif=A.f-S[i].f;
         if (dif<0.0)S[i]=A;
         else
         {
            dif=exp(-dif/Temp);
            if (dif>Lim) S[i]=A;
         }
     }
  }
}
int main (void)
{
   int i, k;
   while (scanf("%d",&nt)==1&&nt)
   {
      maxx=maxy=0;
      minx=miny=(1<<20);
      for (i=0;i<nt;i++)
      {
          scanf("%lf %lf",&T[i].x,&T[i].y);
          if (maxx<T[i].x)maxx=T[i].x;
          if (minx>T[i].x)minx=T[i].x;
          if (maxy<T[i].y)maxy=T[i].y;
          if (miny>T[i].y)miny=T[i].y;
      }
      lx=maxx-minx;
      ly=maxy-miny;
      Temp=sqrt(lx*lx+ly*ly)/3.0;
      ns=5, nc=10;


      Seed();
      while (Temp>EPS)
      {
         Trans();
         Temp=Temp*0.40;
      }
      k=0;
      for (i=1;i<ns;i++)
         if (S[k].f>S[i].f)
            k=i;
      printf ("%.0lf\n", S[k].f);
    }
    return 0;
}



\end{verbatim}
\endgroup
\section{Hệ tọa độ lục giác}
\begingroup
\scriptsize
\begin{verbatim}
//第一种六边形坐标系
#include<stdio.h>
#include<math.h>
#include<string.h>
#include<stdlib.h>
double Dis(double x1, double y1, double x2, double y2)
{
  double dx=x1-x2;
  double dy=y1-y2;
  return sqrt(dx*dx+dy*dy);
}
void Get_KL(double L, double x, double y, int &k, int &l, double &cd)
{
  k=floor((2.0*x)/(3.0*L));
  l=floor((2.0*y)/(sqrt(3.0)*L));
  double d1, d2, x1, y1, x2, y2;
  if ((k+l)&1)
  {
     x1=k*L*1.5;
     y1=(l+1.0)*L*sqrt(3.0)*0.5;
     x2=(k+1.0)*L*1.5;
     y2=l*L*sqrt(3.0)*0.5;
     d1=Dis(x1,y1, x,y);



    d2=Dis(x2,y2, x,y);
    if (d1>d2)
    {
       k++;
       cd=d2;
    }
    else
    {
       l++;
       cd=d1;
    }
  }
  else
  {
     x1=k*L*1.5;
     y1=l*L*sqrt(3.0)*0.5;
     x2=(k+1.0)*L*1.5;
     y2=(l+1.0)*L*sqrt(3.0)*0.5;
     d1=Dis(x1,y1, x,y);
     d2=Dis(x2,y2, x,y);
     if (d1>d2)
     {
        k++,l++;
        cd=d2;
     }
     else cd=d1;
  }
}
int My_Abs(int x)
{
   if (x<0) return -x;
   return x;
}
int main (void)
{
   double L, x1, y1, x2, y2, ans, cd1, cd2;
   int k1, l1, k2, l2;
   while (scanf("%lf %lf %lf %lf %lf",&L,&x1,&y1,&x2,&y2)==5)
   {
      if (L==0.0&&x1==0.0&&y1==0.0&&x2==0.0&&y2==0.0) break;
      Get_KL(L, x1, y1, k1, l1, cd1);
      Get_KL(L, x2, y2, k2, l2, cd2);
      if (k1==k2&&l1==l2) printf("%.3lf\n", Dis(x1,y1, x2,y2));
      else


     {
       ans=cd1+cd2;
       if (My_Abs(k1-k2) > My_Abs(l1-l2))
          ans=ans+sqrt(3.0)*L*My_Abs(k1-k2);
                else ans=ans+sqrt(3.0)*L*My_Abs(k1-k2)+sqrt(3.0)*L*(double)(My_Abs(l1-l2)-
My_Abs(k1-k2))/2.0;
       printf("%.3lf\n", ans);
     }
  }
  return 0;
}

//第二种六边形坐标系
#include <stdio.h>
#include <string.h>
#include <stdlib.h>
#include <math.h>
struct A
{
   int x, y, num;
}a[10001];
const int dec[6][2] = {{-1,1},{-1,0},{0,-1},{1,-1},{1,0},{0,1}};
bool adj(int x1, int y1, int x2, int y2)
{
   if (x1 == x2 && abs(y1-y2) == 1) return true;
   if (y1 == y2 && abs(x1-x2) == 1) return true;
   if (x1 == x2 + 1 && y1 == y2 -1) return true;
   if (x1 == x2 - 1 && y1 == y2 +1) return true;
   return false;
}
bool flag[10001];
int main (void)
{
   int i, j, k, x, u, v, cut, minn, cnt[6];
   memset(cnt, 0, sizeof(cnt));
   a[1].num = 1, cnt[1] = 1;
   a[1].x = a[1].y = 0;
   for (i = 2; i < 10001; i ++)
   {
      k = (int)((3.0+sqrt(12.0*i - 3.0))/6.0+0.0000001);
      if (i == 3*(k-1)*(k-1)+3*(k-1)+1) k --;
      j = i - (3*(k-1)*(k-1)+3*(k-1)+1);
      // 当前的六边形是第 k 层的第 j 个六边形
      if (j == 1) a[i].x = a[i-1].x, a[i].y = a[i-1].y + 1;


      else
      {
         x = (j-1) / k;
         a[i].x = a[i-1].x + dec[x][0], a[i].y = a[i-1].y + dec[x][1];
      }
      memset(flag, false, sizeof(flag));
      x = 12*k-6, cut = 0;
      for (u = i-1, v = 0; u>=1&&v<x; u --, v ++)
         if (adj(a[u].x, a[u].y, a[i].x, a[i].y))
         {
            cut ++;
            flag[a[u].num] = true;
            if (cut == 3) break;
         }
      minn = 10001;
      for (u = 1; u < 6; u ++)
         if ((!flag[u])&&minn > cnt[u])
         {
            minn = cnt[u];
            x = u;
         }
      a[i].num = x;
      cnt[x] ++;
    }
    scanf ("%d", &x);
    while (x --)
    {
       scanf ("%d", &i);
       printf ("%d\n", a[i].num);
    }
    return 0;
}



\end{verbatim}
\endgroup
\section{Dùng đường tròn bán kính cho trước để phủ nhiều điểm nhất}
\begingroup
\scriptsize
\begin{verbatim}
//同半径圆的圆弧表示
#include <stdio.h>
#include <string.h>
#include <stdlib.h>
#include <math.h>



#define PI acos(-1.0)
struct Point
{
   double x, y;
}pt[2005];
double dis[2005][2005];
struct List
{
   double a;
   bool flag;
   int id;
}list[8005];
int cnt;
double Dis(int i, int j)
{
   double dx=pt[i].x-pt[j].x;
   double dy=pt[i].y-pt[j].y;
   return sqrt(dx*dx+dy*dy);
}
int Cmp(const void*p1, const void*p2)
{
   struct List*a1=(struct List*)p1;
   struct List*a2=(struct List*)p2;
   if (a1->a<a2->a)return -1;
   else if (a1->a==a2->a) return a1->id-a2->id;
   else return 1;
}
int main (void)
{
   int n, i, j, ans, num;
   double r, theta, delta, a1, a2;
   while (scanf("%d %lf",&n,&r)==2)
   {
      if (n==0&&r==0.0) break;
      r=r+0.001;
      r=r*2.0;
      for (i=1;i<=n;i++)
         scanf("%lf %lf", &pt[i].x, &pt[i].y);
      for (i=1;i<n;i++)
         for (j=i+1;j<=n;j++)
         {
             dis[i][j]=Dis(i, j);
             dis[j][i]=dis[i][j];
         }


      ans=0;
      for (i=1;i<=n;i++)
      {
         cnt=0;
         for (j=1;j<=n;j++)
            if ((j!=i)&&(dis[i][j]<=r))
            {
               theta=atan2(pt[j].y-pt[i].y, pt[j].x-pt[i].x);
               if (theta<0.0) theta=theta+2.0*PI;
               delta=acos(dis[i][j]/r);
               a1=theta-delta;
               a2=theta+delta;
               list[++cnt].a=a1;
               list[cnt].flag=true;
               list[cnt].id=cnt;
               list[++cnt].a=a2;
               list[cnt].flag=false;
               list[cnt].id=cnt;
            }
         qsort(list+1,cnt,sizeof(struct List),Cmp);
         num=0;
         for (j=1;j<=cnt;j++)
            if (list[j].flag)
            {
               num++;
               if (num>ans) ans=num;
            }
            else num--;
      }
      printf("It is possible to cover %d points.\n", ans+1);
    }
    return 0;
}



\end{verbatim}
\endgroup
\section{Biểu diễn cung của các đường tròn có bán kính khác nhau}
\begingroup
\scriptsize
\begin{verbatim}
intersection_circle_circle(circle[i].center, circle[i].r, circle[j].center, circle[j].r, p1, p2);
             a1= atan2(p1.y-circle[j].center.y, p1.x-circle[j].center.x);
             if (a1<0.0) a1=a1+2.0*PI;
             a2= atan2(p2.y-circle[j].center.y, p2.x-circle[j].center.x);



           if (a2<0.0) a2=a2+2.0*PI;

           if (a1>a2)
           {
              tmp=a1;
              a1=a2;
              a2=tmp;
           }
           mid=(a1+a2)/2.0;
           xtest = circle[j].center.x +circle[j].r*cos(mid);
           ytest = circle[j].center.y +circle[j].r*sin(mid);

           if (!point_in_circle(xtest, ytest, i))
           {
              circle[j].cnt++;
              circle[j].line[circle[j].cnt].s=0;
              circle[j].line[circle[j].cnt].t=a1;
              circle[j].cnt++;
              circle[j].line[circle[j].cnt].s=a2;
              circle[j].line[circle[j].cnt].t=2.0*PI;
           }
           else
           {
              circle[j].cnt++;
              circle[j].line[circle[j].cnt].s=a1;
              circle[j].line[circle[j].cnt].t=a2;
           }



\end{verbatim}
\endgroup
\section{Hợp diện tích hình chữ nhật}
\begingroup
\scriptsize
\begin{verbatim}
#include<stdio.h>
#include<string.h>
#include<stdlib.h>
#include<math.h>
struct Node
{
   int l, r, cnt;
   double cover;
}node[80005];
struct Point



{
   double x;
   double y1, y2;
   int id_y1, id_y2, id_x;
   bool flag;
}pt[20005];
double y[20005];
int total, cnty;
int cmp1(const void*p1, const void*p2)
{
   double*a1=(double*)p1;
   double*a2=(double*)p2;
   if (*a1<*a2) return -1;
   else if (*a1==*a2) return 0;
   else return 1;
}
int cmp2(const void*p1, const void*p2)
{
   struct Point*a1=(struct Point*)p1;
   struct Point*a2=(struct Point*)p2;
   if (a1->x<a2->x) return -1;
   else if (a1->x==a2->x)
   {
      if (a1->id_x<a2->id_x) return -1;
      else if (a1->id_x==a2->id_x) return 0;
      else return 1;
   }
   else return 1;
}
int find(double target)
{
   int head=1, tail=cnty, mid;
   while (head<=tail)
   {
      mid=(head+tail)>>1;
      if (y[mid]==target) return mid;
      else if (y[mid]<target) head=mid+1;
      else tail=mid-1;
   }
   return 0;
}
void Build(int l, int r, int s)
{
   node[s].l=l;


  node[s].r=r;
  node[s].cnt=0;
  node[s].cover=0.0;
  if (l+1<r)
  {
     int mid=(l+r)>>1;
     Build(l,mid,s<<1);
     Build(mid,r,(s<<1)+1);
  }
}
void Update(int s)
{
  if (node[s].cnt>0)
     node[s].cover=y[node[s].r]-y[node[s].l];
  else if(node[s].l+1==node[s].r)
     node[s].cover=0.0;
  else node[s].cover=node[s<<1].cover+node[(s<<1)+1].cover;
}
void Insert(int l, int r, int s)
{
  if (l<=node[s].l&&node[s].r<=r)
  {
     node[s].cnt++;
     Update(s);
     return;
  }
  if (node[s].l+1<node[s].r)
  {
     int mid=(node[s].l+node[s].r)>>1;
     if (l<mid) Insert(l,r,s<<1);
     if (r>mid) Insert(l,r,(s<<1)+1);
     Update(s);
  }
}
void Delete(int l, int r, int s)
{
  if (l<=node[s].l&&node[s].r<=r)
  {
     if (node[s].cnt>0)
        node[s].cnt--;
     Update(s);
     return;
  }
  if (node[s].l+1<node[s].r)


  {
      int mid=(node[s].l+node[s].r)>>1;
      if (l<mid) Delete(l,r,s<<1);
      if (r>mid) Delete(l,r,(s<<1)+1);
      Update(s);
  }
}
int main (void)
{
   int n, i, j, countn=0;
   double ans;
   while (scanf("%d", &n)==1 && n)
   {
      cnty=total=0;
      for (i=1;i<=n;i++)
      {
         total++;
         scanf("%lf %lf", &pt[total].x, &pt[total].y1);
         pt[total].flag=true;
         pt[total].id_x=total;
         y[++cnty]=pt[total].y1;

        total++;
        scanf("%lf %lf", &pt[total].x, &pt[total].y2);
        pt[total].flag=false;
        pt[total].id_x=total;
        y[++cnty]=pt[total].y2;

        pt[total].y1=pt[total-1].y1;
        pt[total-1].y2=pt[total].y2;
      }
      qsort(y+1, cnty, sizeof(double), cmp1);
      j=cnty;
      cnty=1;
      for (i=2;i<=j;i++)
        if (y[i]!=y[i-1])
            y[++cnty]=y[i];

      for (i=1;i<=total;i++)
      {
        pt[i].id_y1=find(pt[i].y1);
        pt[i].id_y2=find(pt[i].y2);
      }
      qsort(pt+1, total, sizeof(struct Point), cmp2);


      ans=0.0;
      Build(1,cnty,1);
      Insert(pt[1].id_y1, pt[1].id_y2, 1);
      for (i=2;i<=total;i++)
      {
         ans=ans+(pt[i].x-pt[i-1].x)*node[1].cover;
         if (pt[i].flag) Insert(pt[i].id_y1, pt[i].id_y2, 1);
         else Delete(pt[i].id_y1, pt[i].id_y2, 1);
      }
      printf("%.0lf\n", ans+1e-10);
    }
    return 0;
}



\end{verbatim}
\endgroup
\section{Hợp chu vi hình chữ nhật}
\begingroup
\scriptsize
\begin{verbatim}
#include <stdio.h>
#include <string.h>
#include <stdlib.h>
struct Point
{
   int x, y;
}plist[10001];
struct Line
{
   int x, b, e, flag;
}llist[10001];
struct Item
{
   int y, id, idx;
}ilist[10001];
struct Node
{
   int l, r, c, m, line;
   bool lf, rf;
}node[40005];
int ys[10001];
int cmp1(const void*p1, const void*p2)
{



  struct Item *a1 = (struct Item*)p1;
  struct Item *a2 = (struct Item*)p2;
  return a1->y - a2->y;
}
int cmp2(const void*p1, const void*p2)
{
   struct Item *a1 = (struct Item*)p1;
   struct Item *a2 = (struct Item*)p2;
   return a1->id - a2->id;
}
int cmp3(const void*p1, const void*p2)
{
   struct Line *a1 = (struct Line*)p1;
   struct Line *a2 = (struct Line*)p2;
   return a1->x - a2->x;
}
void getm(int s)
{
   if (node[s].c > 0)
   {
      node[s].m = ys[node[s].r-1] - ys[node[s].l-1];
      node[s].line = 1;
      node[s].rf = node[s].lf = true;
   }
   else if (node[s].r - node[s].l <= 1)
   {
      node[s].m = node[s].line = 0;
      node[s].rf = node[s].lf = false;
   }
   else
   {
      node[s].m = node[s<<1].m + node[(s<<1)+1].m;
      node[s].line = node[s<<1].line + node[(s<<1)+1].line;
      if (node[s<<1].rf && node[(s<<1)+1].lf) node[s].line --;
      node[s].lf = node[s<<1].lf;
      node[s].rf = node[(s<<1)+1].rf;
   }
}
void build(int l, int r, int s)
{
   node[s].l = l;
   node[s].r = r;
   node[s].c = node[s].m = node[s].line;
   if (node[s].r - node[s].l > 1)


  {
      int mid = (node[s].l + node[s].r)>>1;
      build(l,mid,s<<1);
      build(mid,r,(s<<1)+1);
  }
}
void insert(int l, int r, int s)
{
   if (l <= node[s].l && node[s].r <= r)
   {
      node[s].c ++;
      getm(s);
   }
   if (node[s].r - node[s].l > 1)
   {
      int mid = (node[s].l + node[s].r)>>1;
      if (l < mid) insert(l, r, s<<1);
      if (mid < r) insert(l, r, (s<<1)+1);
      getm(s);
   }
}
void delet(int l, int r, int s)
{
   if (l <= node[s].l && node[s].r <= r)
   {
      node[s].c --;
      getm(s);
   }
   if (node[s].r - node[s].l > 1)
   {
      int mid = (node[s].l + node[s].r)>>1;
      if (l < mid) delet(l, r, s<<1);
      if (mid < r) delet(l, r, (s<<1)+1);
      getm(s);
   }
}
int main (void)
{
   int n, i, j, l, r, x1, y1, x2, y2, tot, p, ans;
   while (scanf ("%d", &n) == 1 && n)
   {
      for (i = 0; i < n; i ++)
      {
         scanf ("%d %d %d %d", &x1, &y1, &x2, &y2);


  l = 2*i;
  r = l + 1;

  plist[l].x = x1;
  plist[l].y = y1;
  plist[r].x = x2;
  plist[r].y = y2;

  ilist[l].y = y1;
  ilist[l].id = l;
  ilist[r].y = y2;
  ilist[r].id = r;
}
tot = 2*n;
qsort(ilist, tot, sizeof(struct Item), cmp1);
ys[0] = ilist[0].y;
ilist[0].idx = 0;
j = 0;
for (i = 1; i < tot; i ++)
{
   if (ilist[i].y != ilist[i-1].y)
   {
      j ++;
      ys[j] = ilist[i].y;
   }
   ilist[i].idx = j;
}
p = j + 1;
qsort(ilist, tot, sizeof(struct Item), cmp2);
for (i = 0; i < n; i ++)
{
   l = 2*i;
   r = l + 1;
   llist[l].x = plist[l].x;
   llist[l].b = ilist[l].idx;
   llist[l].e = ilist[r].idx;
   llist[l].flag = 1;

  llist[r].x = plist[r].x;
  llist[r].b = ilist[l].idx;
  llist[r].e = ilist[r].idx;
  llist[r].flag = 0;
}
qsort(llist, tot, sizeof(struct Line), cmp3);


      build(1,p,1);
      insert(llist[0].b+1, llist[0].e+1,1);
      int now_m = node[1].m, now_line = node[1].line;
      ans = now_m;
      for (i = 1; i < tot; i ++)
      {
         if (llist[i].flag) insert(llist[i].b+1, llist[i].e+1, 1);
         else delet(llist[i].b+1, llist[i].e+1, 1);
         ans += (abs(node[1].m - now_m) + 2*(llist[i].x - llist[i-1].x)*now_line);
         now_m = node[1].m;
         now_line = node[1].line;
      }
      printf ("%d\n", ans);
    }
    return 0;
}



\end{verbatim}
\endgroup
\section{Cặp đường tròn gần nhất}
\begingroup
\scriptsize
\begin{verbatim}
#include<iostream>
#include<stdlib.h>
#include<string.h>
#include<set>
#include <math.h>
using namespace std;
set <int>tree;
set <int>::iterator iter;
struct Point
{
   double x;
   int id, flag;
}p1[100001], p2[100001];
int tot1, tot2;
struct Q
{
   double x,y, r;
}q[50001];
int cmp(const void*p1, const void*p2)
{
   struct Point*a1=(struct Point*)p1;



  struct Point*a2=(struct Point*)p2;
  if (a1->x<a2->x) return -1;
  else if (a1->x==a2->x) return a2->flag-a1->flag;
  else return 1;
}
int cmp1(const void*p1, const void*p2)
{
   struct Q*a1=(struct Q*)p1;
   struct Q*a2=(struct Q*)p2;
   if (a1->y<a2->y)return -1;
   else if (a1->y==a2->y)return 0;
   else return 1;
}
double dis(double x1, double y1, double x2, double y2)
{
   return sqrt((x1-x2)*(x1-x2)+(y1-y2)*(y1-y2));
}
bool judge(int i, int j, double d)
{
   if (dis(q[i].x, q[i].y, q[j].x, q[j].y)<=q[i].r+q[j].r+2.0*d)
      return true;
   return false;
}
bool insert(int v,double d)
{
   iter = tree.insert(v).first;
   if (iter != tree.begin())
   {
      if (judge(v, *--iter,d))
      {
          return true;
      }
      ++iter;
   }
   if (++iter != tree.end())
   {
      if (judge(v, *iter,d))
      {
          return true;
      }
   }
   return false;
}
bool remove(int v,double d)


{
    iter = tree.find(v);

    if (iter != tree.begin() && iter != --tree.end())
    {
       int a = *--iter;
       ++iter;
       int b = *++iter;
       if (judge(a, b,d))
       {
           return true;
       }
    }
    tree.erase(v);
    return false;
}
bool check(double d)
{
  int i=1, j=1;
  while (i<=tot1&&j<=tot2)
  {
     if (p1[i].x-d<=p2[j].x+d)
     {
        if (insert(p1[i++].id, d))
           return true;
     }
     else
     {
        if (remove(p2[j++].id, d))
           return true;
     }

    }
    while (i<=tot1)
    {
       if (insert(p1[i++].id, d))
          return true;
    }
    while (j<=tot2)
    {
       if (remove(p2[j++].id, d))
          return true;
    }
    return false;


}
int main (void)
{
   int cases, n, i;
   scanf("%d",&cases);
   while (cases--)
   {
      scanf("%d",&n);
      tot1=tot2=0;
      for (i=1;i<=n;i++)
         scanf("%lf %lf %lf",&q[i].x,&q[i].y, &q[i].r);
      qsort(q+1,n,sizeof(struct Q),cmp1);
      for (i=1;i<=n;i++)
      {
         tot1++;
         p1[tot1].x=q[i].x-q[i].r;
         p1[tot1].id=i;
         p1[tot1].flag=1;

         tot2++;
         p2[tot2].x=q[i].x+q[i].r;
         p2[tot2].id=i;
         p2[tot2].flag=-1;
      }
      qsort(p1+1,tot1,sizeof(struct Point),cmp);
      qsort(p2+1,tot2,sizeof(struct Point),cmp);

      double head=0.0, tail=dis(q[1].x,q[1].y,q[2].x,q[2].y)+1.0, mid;
      while (tail-head>1e-8)
      {
         tree.clear();
         mid=(head+tail)/2.0;
         if (check(mid))
         {
            tail=mid;
         }
         else head=mid;
      }
      printf ("%.6lf\n",2.0*head);
    }
    return 0;
}





\end{verbatim}
\endgroup
\section{Tính diện tích giao của hai đường tròn}
\begingroup
\scriptsize
\begin{verbatim}
double area_of_overlap(point c1, double r1, point c2, double r2)
{
  double a = distance(c1, c2), b = r1, c = r2;
    double cta1 = acos((a * a + b * b - c * c) / 2 / (a * b)),
        cta2 = acos((a * a + c * c - b * b) / 2 / (a * c));
  double s1 = r1*r1*cta1 - r1*r1*sin(cta1)*(a * a + b * b - c * c) / 2 / (a * b);
  double s2 = r2*r2*cta2 - r2*r2*sin(cta2)*(a * a + c * c - b * b) / 2 / (a * c);
  return s1 + s2;
}





\end{verbatim}
\endgroup
\end{document}
